\documentclass[11pt,a4paper]{article}

\usepackage[utf8]{inputenc}
\usepackage[T1]{fontenc}
\usepackage[english]{babel}
\usepackage{amsmath,amssymb,amsthm}
\usepackage{mathtools}
\usepackage{geometry}
\geometry{margin=2.5cm}
\usepackage{booktabs}
\usepackage{longtable}
\usepackage{enumitem}
\usepackage[hidelinks]{hyperref}

\theoremstyle{plain}
\newtheorem{theorem}{Theorem}
\newtheorem{proposition}[theorem]{Proposition}
\newtheorem{lemma}[theorem]{Lemma}
\newtheorem{corollary}[theorem]{Corollary}
\theoremstyle{definition}
\newtheorem{postulate}{Postulate}
\newtheorem{result}{Spectral Result}
\theoremstyle{remark}
\newtheorem{remark}[theorem]{Remark}

\newcommand{\Ok}{\Omega_K}
\newcommand{\OL}{\Omega_\Lambda}
\newcommand{\Om}{\Omega_m}
\newcommand{\ag}{a_\gamma}
\newcommand{\CQ}{C_Q}
\newcommand{\Lb}{\mathcal{L}}
\newcommand{\Mpl}{M_{\mathrm{Pl}}}
\newcommand{\dtheta}{\delta\theta}
\newcommand{\Rtop}{\mathcal{R}}
\newcommand{\avg}[1]{\langle #1 \rangle}
\DeclareMathOperator{\Vol}{Vol}
\DeclareMathOperator{\tr}{tr}

\title{\bfseries A Conformal-Anomaly Resolution of the Cosmological Constant Problem\\[0.4em]
\large The Vacuum Energy, Its Value, the Coincidence, and a Parameter-Free
Curvature Fingerprint --- with a Complete Account of What Is Proven and What Is Conditional}
\author{Lukas R\"ohl\thanks{Independent Researcher, Kelmis (Belgium).}}
\date{}

\begin{document}
\maketitle

\begin{abstract}
\noindent
The cosmological constant problem has three faces: the \emph{magnitude} problem (why the
observed vacuum energy is $\sim 10^{-120}$ in Planck units rather than of order unity), the
\emph{value} problem (what fixes the residual scale once the radiative contributions are
removed), and the \emph{coincidence} problem (why the dark-energy and matter densities are
comparable now). This paper assembles a single, internally consistent framework that addresses
all three under explicit, stated assumptions --- notably the physical identifications A5 and P5
made precise below --- and, as a by-product, yields one sharp falsifiable prediction. The magnitude problem
is handled by Kaloper--Padilla vacuum-energy sequestering, which removes the radiative vacuum
energy from the gravitating sector. The value problem is addressed by a self-consistency
result (Proposition~\ref{thm:kp}): the residual cosmological constant is a fixed point of
the four-volume averaging map $\Lambda = \tfrac14\avg{R}_{\mathrm{ren}}$, equivalent to the
trace-zero condition $\avg{T}_{\mathrm{ren}} = 0$, and is therefore \emph{not} a free
parameter; the uniqueness of that fixed point rests on a Banach contraction whose constant is so
far established numerically. The coincidence problem is addressed by structure-formation (halo-weighted)
averaging following Lombriser, returning $\OL \approx 0.704$ conditional on an anthropic-like
prior; the framework's own budget ladder instead points to $\OL = 4\ag = 0.6889$,
coinciding with the robust Planck+BAO value ($\Om = 0.3111$) to four decimals --- an empirical
coincidence, not a proof: the absolute value rests on the open holographic-saturation conjecture
P5 (Section~\ref{sec:ladder}). Threading through all three is a relation --- parameter-free as a
ratio of two scale-free invariants, though its identification with the \emph{observable} curvature
rests on the physical postulate A5 (below) --- between the residual cosmological constant and the
spatial curvature,
\[
\frac{|\Ok|}{\OL} \;=\; \frac{\ag}{8\pi^2} \;=\; \frac{31}{1440\pi^2} \;\approx\; 2.181\times 10^{-3},
\qquad
\Lb \equiv \frac{\OL}{|\Ok|} = \frac{8\pi^2}{\ag} \approx 458.4,
\]
whose two inputs --- the photon type-A conformal anomaly coefficient $\ag = 31/180$ and the
$Q$-curvature charge $\CQ = \int_{D^4} Q_4 = 8\pi^2$ of the hemisphere --- contain no scale, no
Newton constant, and no Planck mass. We give the rigorous spectral backbone in full: twelve
proven spectral results (R1--R12), the Burghelea--Friedlander--Kappeler gluing identity, the
cap-variational saddle and its stability in every sector, the higher-loop convergence bound,
and an independent recomputation (M1) with exact symbolic and numerical methods. Every step
carries an explicit status --- \emph{proven}, \emph{literature-supported}, \emph{calibrated},
or \emph{open}. After decomposition, the framework reduces to a single genuinely chosen
element --- the data-anchored reference state A5c --- plus the embedding into the full
infinite-dimensional no-boundary path integral (the two-dimensional cap-plus-scalar embedding
$B1'$ being complete, $n_{\mathrm{complex}}=0$ by polynomial reduction); its
topological core (the ratio A5a, the self-consistency fixed point A2, cap dominance B1 on the
reduced action) is theorem-level, and the open sign is forced by the positive photon anomaly coefficient $\ag>0$ \emph{within the Hartle--Hawking branch} (the branch selection itself a postulate).
The prediction is
sharply falsifiable: a robust $\Ok$ outside $\pm 3\sigma$ of $\ag\OL/(8\pi^2)$, or a robust
detection of a closed universe, would refute the framework as stated.
\end{abstract}

\tableofcontents
\bigskip

\section{The cosmological constant problem and what a resolution must deliver}
\label{sec:intro}

The cosmological constant $\Lambda$ is the sharpest quantitative conflict between quantum field
theory and general relativity: naive estimates of the quantum-vacuum zero-point energy exceed the
observed dark-energy density by some $120$ orders of magnitude. The problem is not one question but three, and any framework claiming a
\emph{resolution} must say something about each.

\paragraph{(i) The magnitude problem.} Why is the gravitating vacuum energy
$\rho_\Lambda \sim 10^{-120}\,\Mpl^4$ rather than $O(\Mpl^4)$? Every massive field contributes
a quartically divergent vacuum energy; the cancellation to $120$ digits is the heart of the
old cosmological constant problem.

\paragraph{(ii) The value problem.} Suppose the radiative contributions are removed by some
mechanism. What then fixes the residual, nonzero scale? A mechanism that sets $\Lambda$ to zero
is no better than one that leaves it free: the observed value is small but not zero.

\paragraph{(iii) The coincidence problem.} Why are $\OL \approx 0.69$ and $\Om \approx 0.31$ of
the same order \emph{now}, when their ratio scales as $a^3$ and has therefore been wildly
different at almost every other epoch?

\paragraph{The strategy of this paper.} We address (i) by Kaloper--Padilla sequestering
(Section~\ref{sec:sequester}), which removes the radiative vacuum energy from the source of
Einstein's equations. We address (ii) by a self-consistency result
(Section~\ref{sec:theorem64}): the residual $\Lambda$ is a fixed point of a
four-volume averaging map, hence dynamically determined rather than free. We address (iii) by
structure-formation averaging (Section~\ref{sec:coincidence}), which returns the observed
$\OL \approx 0.704$ in Lombriser's halo-weighted prescription, conditional on an anthropic-like
prior; the framework's own --- conjectural --- budget ladder instead points
to $\OL = 4\ag = 0.6889$, coinciding with the robust Planck+BAO datum to four decimals --- an
empirical coincidence, not a proof (Section~\ref{sec:ladder}).
The framework's sharpest --- and only sharply falsifiable --- output is a relation, parameter-free as a ratio of two scale-free invariants,
between the residual cosmological constant and
the present spatial curvature,
\begin{equation}
\label{eq:central}
\boxed{\;\frac{|\Ok|}{\OL} = \frac{\ag}{8\pi^2} = \frac{31}{1440\pi^2} \approx 2.181\times 10^{-3}\;}
\end{equation}
with $\ag = 31/180$ the photon type-A conformal anomaly coefficient and $8\pi^2$ the
$Q$-curvature charge of the four-hemisphere $D^4 \subset S^4$. In the present epoch the ladder-consistent budget
$\Om = 0.3096$ gives, by the Friedmann sum rule, $|\Ok| \approx 1.50\times 10^{-3}$ and
$\OL = 4\ag = 0.6889$ (Section~\ref{sec:ladder}).

\paragraph{What is claimed, and what is not.} The account rests on three disclaimers, stated plainly at the outset.
\begin{enumerate}[leftmargin=*,nosep]
\item The relation~\eqref{eq:central} is a statement about a \emph{ratio}. The \emph{absolute}
dark-energy density is not a \emph{sharp} parameter-free output of the topology: the KP
self-consistency route fixes $\Lambda$ \emph{structurally} (not free, sign negative, the bulk
vacuum cancels) and delivers a \emph{universal} fixed point $\Lambda^*$ (independent of
$\varphi_0$ and $\Omega_m$), but its value is amplitude-set and the absolute scale is left open
because $\varphi_0$ stays free (Section~\ref{sec:theorem64}; the earlier $\langle R\rangle/4$
sharpness proxy was a floored-geometry artefact). The sharp value is then either supplied by
structure-formation (halo) averaging under one additional assumption
(Section~\ref{sec:coincidence}), bracketed prior-free by the causal closure
(Section~\ref{sec:measure}), or --- conjecturally --- fixed topologically by the geometric budget
ladder $\OL = 4\ag$ (Section~\ref{sec:ladder}), whose lift prefactor is computed ($=A/G$, the
Bekenstein--Hawking horizon degree-of-freedom count) and whose residual --- the holographic
Cohen--Kaplan--Nelson saturation P5 --- is localized but not proven. This residual is
confined to the \emph{absolute} value: the ratio $\Ok/\OL=\ag/8\pi^2$ is scale-invariant, so the
falsifiable Euclid prediction stands independently of it.
\item The Euclidean conformal and topological core is theorem-level --- including the
topological ratio A5a (Cauchy + Chern--Weil + APS), the KP self-consistency fixed point A2
(M5 Banach contraction), and cap-saddle dominance B1 on the reduced action ($n_{\mathrm{complex}}
= 0$, analytic). The bridge to the observed Lorentzian Friedmann curvature rests on a single
remaining physical postulate, the identification A5b/A5c (Section~\ref{sec:main}), decided by
observation; the two-dimensional cap-plus-scalar embedding $B1'$ is itself now complete ($n_{\mathrm{complex}}=0$ by a degree-$8$ polynomial reduction, Section~\ref{sec:found}), and that embedding into the full infinite-dimensional no-boundary path integral is assembled semiclassically (Section~\ref{sec:found}): the inhomogeneous one-loop determinant carries the anomaly ($\zeta(0)=-4\ag$) and the Type-A coefficient is Adler--Bardeen one-loop exact, and the prediction $\ag/(8\pi^2)$ is all-orders stable as a ratio of $\hbar$-independent invariants (Section~\ref{sec:found}); the only residue is the absolute Borel summability of the matter series, decoupled from the prediction.
\item The \emph{sign} (open versus closed) is forced by the positive photon anomaly coefficient
$\ag>0$: the cap action has a unique real saddle at $\dtheta^*>0$ whose location is independent
of the Hartle--Hawking/Vilenkin weighting (Section~\ref{sec:sign}, Section~\ref{sec:found}). The
separate question of Lorentzian fluctuation stability (Feldbrugge--Lehners--Turok) is a
foundational input of all no-boundary cosmology, not specific to this framework.
\end{enumerate}

\paragraph{A historical analogy.} A clean, parameter-free relation justified at the level of a
topological/anomaly identity, with a physically motivated but not yet rigorously derived bridge
to the observable, has a precedent. Bekenstein's 1972 conjecture that black-hole entropy is
proportional to horizon area \cite{Bekenstein1972} was a thermodynamic argument; its
microscopic derivation arrived twenty-three years later with Strominger and Vafa
\cite{Strominger1996}. The present framework is offered in the same spirit.

\paragraph{Outline.} Sections~\ref{sec:anomaly}--\ref{sec:spectral} assemble the three rigorous
ingredients: the photon anomaly coefficient $\ag$, the $Q$-curvature charge $\CQ$, and the
spectral backbone (twelve results R1--R12 and the BFK identity). Sections~\ref{sec:sequester}
and~\ref{sec:theorem64} treat the magnitude and value problems (Kaloper--Padilla sequestering
and the self-consistency result). Section~\ref{sec:main} derives
relation~\eqref{eq:central} --- parameter-free given the identification A5 --- in eight steps and isolates the residual reference choice A5c;
Section~\ref{sec:cap} gives the cap-variational mechanism, its stability, the boundary
refinement, and the prediction hierarchy; Section~\ref{sec:uniqueness} establishes uniqueness,
within stated criteria, against $26$ alternative mechanisms. Section~\ref{sec:coincidence} addresses the coincidence;
Section~\ref{sec:sign} establishes the sign within the Hartle--Hawking branch; Section~\ref{sec:psf} gives the physical
interpretation; Section~\ref{sec:corrections} the higher-order corrections;
Section~\ref{sec:m1} the independent computational verification (M1);
Section~\ref{sec:obs} the observational status; and Section~\ref{sec:status} the consolidated
status ledger and falsification criteria.

\section{The conformal anomaly and the photon coefficient}
\label{sec:anomaly}

\subsection{Trace anomaly and the type-A coefficient}
A classically conformally invariant field has a traceless stress tensor, $T^\mu{}_\mu = 0$.
Quantization breaks this; in four curved dimensions the renormalized trace acquires the
universal form \cite{DS1993,Duff1994}
\begin{equation}
\label{eq:anomaly}
\avg{T^\mu{}_\mu} = \frac{1}{16\pi^2}\Big( c\, C_{\mu\nu\rho\sigma}C^{\mu\nu\rho\sigma}
  \;-\; a\, E_4 \;+\; b\,\Box R \Big),
\end{equation}
where $C$ is the Weyl tensor, $E_4 = R_{\mu\nu\rho\sigma}R^{\mu\nu\rho\sigma}
- 4 R_{\mu\nu}R^{\mu\nu} + R^2$ is the Gauss--Bonnet (Euler) density, and $a$, $c$ are the
type-A and type-B central charges. The $b\,\Box R$ term is scheme-dependent and integrates to
zero on a closed manifold. The Deser--Schwimmer theorem \cite{DS1993} establishes that this
classification is complete and geometric.

\subsection{Why only the photon survives in the infrared}
On the present-day de Sitter-like background the Hubble scale $H_0$ lies far below the mass of
every massive field. By the Komargodski--Schwimmer monotonicity of the $a$-coefficient under
renormalization-group flow \cite{KomargodskiSchwimmer2011}, massive fields decouple from the
infrared conformal sector; only the genuinely massless gauge field --- the photon --- still
contributes to the type-A anomaly at cosmological scales. This is the content of Axiom A3
(photon-uniqueness), placed on a quantitative footing in Section~\ref{sec:uniqueness} by the
exhaustion of $26$ standard-Lagrangian alternatives.

\subsection{The value $\ag = 31/180$}
The photon type-A coefficient is a textbook result of the DeWitt--Schwinger heat-kernel
expansion for a Maxwell field \cite{Birrell1982,ChristensenDuff1980,Duff1994}:
\begin{equation}
\label{eq:agamma}
\ag = \frac{31}{180}.
\end{equation}
We \emph{derive} Eq.~\eqref{eq:agamma} from the Gilkey heat-kernel coefficient $a_4$ in
Section~\ref{sec:heat}: the Maxwell combination (Hodge Laplacian on $1$-forms minus two
Faddeev--Popov ghosts) gives $\zeta(0;\mathrm{Maxwell};S^4) = -31/45$, hence $\ag = 31/180$, with
the conformal-scalar value $1/360$ serving as an independent check. The same coefficient governs
the closed-manifold zeta value
$\zeta(0;\text{Maxwell};S^4) = -4\ag = -31/45$ (Spectral Result R1 below), since the Weyl term
vanishes on the conformally flat $S^4$ and $\int_{S^4} E_4 = 64\pi^2$ gives
$\int_{S^4}\avg{T^\mu{}_\mu} = 4\ag$.

\section{Heat-kernel coefficients}
\label{sec:heat}
The DeWitt--Schwinger expansion of a second-order elliptic operator $D = -\Box + E$ has
$K(x,x;s) = (4\pi s)^{-n/2}[\mathbf{1} + s\,a_1 + s^2 a_2 + \dots]$. In $n=4$ the anomaly-relevant
coefficient is
\[
a_2 = \tfrac{1}{180}(R_{\mu\nu\rho\sigma}^2 - R_{\mu\nu}^2)\mathbf{1} + \tfrac{1}{72}R^2\mathbf{1}
 - \tfrac{1}{30}\Box R\,\mathbf{1} + \tfrac16\Box(\tfrac{R}{6}\mathbf{1} - E)
 + \tfrac12(\tfrac{R}{6}\mathbf{1} - E)^2 + \tfrac{1}{12}W_{\mu\nu}W^{\mu\nu}.
\]
For the Maxwell vector ($E^\nu{}_\mu = -R^\nu{}_\mu$, $\mathrm{tr}\,\mathbf{1} = 4$) plus the
Faddeev--Popov ghost ($E=0$, contributing with a $-2$), the trace gives
$\mathrm{tr}\,a_2 = \tfrac{1}{16\pi^2}[\ag E_4 - c_\gamma W^2 + b\,\Box R]$ with $\ag = 31/180$,
$c_\gamma = 1/10$, and $b$ a removable counterterm.

\paragraph{Explicit evaluation of $\ag$ on $S^4$.} The coefficient $\ag$ is not assumed; it
follows from Gilkey's integrated $a_4$. Dropping total derivatives on a closed manifold,
\[
\int a_4 = \frac{1}{(4\pi)^2\,360}\int\!\sqrt g\;\mathrm{tr}\Big\{60\,R\,E + 180\,E^2
+ 30\,\Omega_{\mu\nu}\Omega^{\mu\nu} + \mathbf{1}\big(5R^2 - 2R_{\mu\nu}^2 + 2R_{\mu\nu\rho\sigma}^2\big)\Big\},
\]
with $\Omega_{\mu\nu}$ the bundle curvature of $D$. On $S^4$ one has $R_{\mu\nu}^2 = R^2/4$,
$R_{\mu\nu\rho\sigma}^2 = R^2/6$, $W^2 = 0$, and $\int E_4 = 64\pi^2$, so only the type-A (Euler)
part survives and the integrated conformal anomaly equals $\int\langle T^\mu{}_\mu\rangle =
-4\ag$ for a single field. Three operators are needed:
\begin{itemize}[leftmargin=*,nosep]
\item \emph{Conformal scalar} ($D=-\Box + \tfrac{R}{6}$, $E=-\tfrac{R}{6}$, $\dim 1$, $\Omega=0$):
the bracket reduces to $-\tfrac16 R^2$, giving $\int\langle T\rangle = -\tfrac{1}{90}$ and
$a_{\mathrm{scalar}} = \tfrac{1}{360}$ --- the textbook value, a check on the whole computation.
\item \emph{Hodge Laplacian on $1$-forms} ($E = -\mathrm{Ric}$, $\mathrm{tr}\,E = -R$,
$\mathrm{tr}\,E^2 = R^2/4$, $\dim 4$, $\mathrm{tr}\,\Omega_{\mu\nu}\Omega^{\mu\nu} =
-R_{\mu\nu\rho\sigma}^2 = -R^2/6$): $\int\langle T\rangle_{\Delta_1} = -\tfrac{2}{45}$.
\item \emph{Minimal scalar ghost} ($E=0$, $\dim 1$): $\int\langle T\rangle_{\Delta_0} =
\tfrac{29}{90}$.
\end{itemize}
The photon partition function is $\ln Z = -\tfrac12\ln\det\Delta_1 + \ln\det\Delta_0^{\mathrm{ghost}}$,
so its anomaly is the combination $\Delta_1 - 2\,\Delta_0$:
\[
\zeta(0;\mathrm{Maxwell};S^4) = \int\langle T\rangle_{\Delta_1} - 2\int\langle T\rangle_{\Delta_0}
= -\tfrac{2}{45} - 2\cdot\tfrac{29}{90} = -\tfrac{31}{45},
\qquad \ag = -\tfrac14\Big(-\tfrac{31}{45}\Big) = \boxed{\tfrac{31}{180}.}
\]
This reproduces Spectral Result R1 and fixes the central input of the framework from the
heat-kernel coefficient alone, so $\ag$ is derived here rather than adopted as an external citation.

\section{The $Q$-curvature charge of the hemisphere}
\label{sec:qcharge}

\subsection{The $Q$-curvature and the Paneitz operator}
Branson's fourth-order $Q$-curvature in four dimensions is
\begin{equation}
Q_4 = -\tfrac{1}{6}\Box R - \tfrac{1}{2} R_{\mu\nu}R^{\mu\nu} + \tfrac{1}{6} R^2 ,
\end{equation}
the natural curvature scalar transforming covariantly under the Paneitz operator $P_4$ via
$P_4\,\sigma + Q_4 = Q_4'\, e^{4\sigma}$ under a Weyl rescaling $g \to e^{2\sigma} g$. The
Chang--Yang theory of extremal functional determinants \cite{ChangYang1995,Branson1995} singles
out $\int Q_4$ as the conformally natural total curvature: only $\int Q_4$ (not $\int E_4$)
transforms covariantly under $P_4$, which is why $Q_4$, and hence the charge $\CQ$, is the
correct denominator of the topological quotient~\eqref{eq:central}. On any closed Einstein
four-manifold $\int Q_4 = \tfrac14\int E_4$ exactly (Spectral Result R12).

\subsection{The hemisphere value $\CQ = 8\pi^2$, with an elementary proof}
\begin{lemma}[Hemisphere $Q$-charge]
\label{lem:qcharge}
On the round unit hemisphere $D^4 = \{0 \le \theta \le \pi/2\} \subset S^4$,
$\CQ \equiv \int_{D^4} Q_4\,\sqrt{g}\,d^4x = 8\pi^2$.
\end{lemma}
\begin{proof}
On the unit $S^4$ (an Einstein manifold) $R = 12$, $R_{\mu\nu} = 3 g_{\mu\nu}$, $\Box R = 0$,
hence $Q_4 = \tfrac16(R^2 - 3 R_{\mu\nu}R^{\mu\nu}) = \tfrac16(144 - 108) = 6$. With the metric
volume element $\sqrt{g} = \sin^3\theta$ and $\Vol(S^3) = 2\pi^2$,
\[
\int_0^{\pi/2}\sin^3\theta\,d\theta = \Big[-\cos\theta + \tfrac13\cos^3\theta\Big]_0^{\pi/2}
 = \tfrac23,
\qquad
\int_{D^4} Q_4\,\sqrt{g}\,d^4x = 6 \cdot 2\pi^2 \cdot \tfrac23 = 8\pi^2 . \qedhere
\]
\end{proof}
\noindent
This elementary three-step evaluation has been confirmed by exact symbolic computation
(Section~\ref{sec:m1}, M1-Q). The constant $8\pi^2$ is the exact, scale-free $Q$-curvature
charge of the hemisphere; it is a \emph{proven} input. The cap-saddle shift
$\dtheta^* \approx 0.164$ introduced in Section~\ref{sec:cap} is a position in the wave-function
argument, \emph{not} a deformation of this integration domain; the charge is always evaluated
on the standard hemisphere, so the leading prediction $|\Ok|/\OL = \ag/(8\pi^2)$ carries no
$\dtheta^*$ correction. The cap-saddle enters only at the refinement level L1, where it
replaces the numerator $\ag$ while the denominator $8\pi^2$ stays fixed.

\section{Spectral foundations: BFK decomposition and the twelve results}
\label{sec:spectral}

The bridge between the closed $S^4$ and its hemisphere $D^4$ is the Burghelea--Friedlander--
Kappeler (BFK) gluing formula for functional determinants \cite{BFK1992}, which relates the
determinant on a closed manifold to those on two pieces glued along a hypersurface plus the
Dirichlet-to-Neumann (DtN) operator $\Lambda_{S^3}$ on the equatorial $S^3$. For the
zeta-regularized value at zero,
\begin{equation}
\label{eq:bfk}
\zeta(0;\Delta_{S^4}) = 2\,\zeta(0;\Delta_{D^4}^{\mathrm{Dir}}) - \zeta(0;\Lambda_{S^3}).
\end{equation}

\subsection{The twelve spectral results R1--R12}
The framework rests on twelve proven spectral results, all independently verified (several to
high numerical precision). They are the rigorous backbone of the construction; we state them
compactly.
\begin{description}[leftmargin=2.6em,nosep,itemsep=2pt]
\item[R1.] $\zeta(0;\mathrm{Maxwell};S^4) = -4\ag = -31/45$. Two derivations agree: the
integrated trace-anomaly route ($-\ag\int_{S^4}E_4/16\pi^2 = -4\ag$) and the spectral
decomposition of the Hodge Laplacian on transverse $1$-forms with Faddeev--Popov ghost
subtraction \cite{ChristensenDuff1980,Vassilevich2003}. Each unit of Euler characteristic
contributes $-2\ag$; $\chi(S^4)=2$ gives $-4\ag$.
\item[R2.] $\zeta(0;\mathrm{DtN};\mathrm{scalar};S^3) = -1$. The scalar DtN operator on the
equatorial $S^3$ of $S^4$ (harmonic extension into the cap) has spectrum
$\lambda_\ell = \ell(\ell+2)/(\ell+1) = (\ell+1)-1/(\ell+1)$ ($\ell\ge1$) with multiplicity
$(\ell+1)^2$ (the constant $\ell=0$ mode is the kernel); with $n=\ell+1$ the leading term
$\sum_{n\ge2}n^{2-s}$ continues to $\zeta_R(s-2)-1$ and the subleading $(1-1/n^2)^{-s}$
corrections all carry a factor $s$ and vanish at $s=0$, giving $\zeta(0)=\zeta_R(-2)-1 = -1$
\cite{Vassilevich2003}. (The flat-ball value $\lambda_\ell=\ell$ does \emph{not} apply to the
$S^4$ cap; it would give $-2/3$.)
\item[R3.] $\zeta(0;\mathrm{DtN};\text{1-form};S^3) = 1$.
\item[R4.] $\zeta(0;\mathrm{DtN};\mathrm{Maxwell};S^3) = 3 = \dim S^3$ (gauge-fixed Maxwell DtN).
\item[R5.] $\zeta(0;\mathrm{Maxwell};S^3) = 1$ (the photon-specific equatorial identity).
\item[R6.] BFK identity verified: with the closed value R1 and the two consistent bookkeepings
of Eq.~\eqref{eq:bfk} (Section~\ref{sec:two-books}), $2\zeta(0;D^4)-\zeta(0;\mathrm{DtN}) = -31/45$.
\item[R7.] Photon Casimir energy on $S^3$: $E_{\mathrm{Cas}} = 11/120$.
\item[R8.] Functional determinant ${\det}'\Lambda_{S^3}^{\mathrm{Maxwell}} = 1/(2\pi^3)$.
\item[R9.] $\eta(\Lambda_{S^3}^{\text{transverse 1-form DtN}}) = \zeta(0) = 1$, conformally invariant
(the one-form constituent of R8; the full Maxwell DtN has $\zeta(0)=3$). (Distinct
from the de~Rham APS eta $\eta_{\mathrm{APS}}(S^3)=0$, which enters the topological quantization
of Section~\ref{sec:main}; they are different operators on the same $S^3$.)
\item[R10.] A Robin condition $(\partial_n + \alpha)u = 0$ at the deformed equator shifts the DtN
spectrum by $\alpha = \tan\ag \approx \ag$ --- the analytic origin of the cap-saddle drive.
\item[R11.] Under the Robin shift, $\zeta(0;\Lambda_{\mathrm{Robin}})$ shifts linearly in $\alpha$.
\item[R12.] $\int_{D^4} Q_4 / \int_{D^4} E_4 = 1/4$ exactly, fixing the $Q$-vs-$E$ normalization.
\end{description}

\subsection{The two consistent BFK bookkeepings}
\label{sec:two-books}
The hemisphere zeta value carries a nontrivial spectral weight, and the BFK
identity~\eqref{eq:bfk} admits two internally consistent bookkeepings that differ only in how
the gauge/ghost modes are assigned between the hemisphere and the DtN operator. Both return the
closed value $\zeta(0;S^4)=-31/45$.
\begin{align}
\text{(a) scalar-DtN split:}\quad
&\zeta(0;D^4) = -2\ag - \tfrac12 = -\tfrac{38}{45},
&&2\big(-\tfrac{38}{45}\big) - (-1) = -\tfrac{31}{45}\;\checkmark \\
\text{(b) gauge-fixed Maxwell split:}\quad
&\zeta(0;D^4) = \tfrac{52}{45},
&&2\big(\tfrac{52}{45}\big) - 3 = -\tfrac{31}{45}\;\checkmark
\end{align}
The equivalence of the two bookkeepings is itself a check on the spectral arithmetic; the
gauge-fixed Maxwell DtN value $3$ (R4) and the symmetric hemisphere value $52/45$ are used in
the boundary analysis of Section~\ref{sec:cap}. These spectral results are \emph{proven}; the
closed value $\zeta(0;\mathrm{Maxwell};S^4) = -31/45$ is itself \emph{derived} from the Gilkey
heat-kernel coefficient in Section~\ref{sec:heat} (cross-checked against Christensen--Duff 1980;
Esposito 1995, in the conventions of \cite{Vassilevich2003}). The
photon-specific value $\zeta(0;\mathrm{Maxwell};S^3)=1$ (R5), together with the DtN data,
underlies the cap-variational restoring force of Section~\ref{sec:cap}.

\section{Spectral calculations: R1--R12}
\label{sec:spectraldetail}
This section collects the explicit computations behind the twelve spectral results R1--R12
stated above.

\paragraph{R1: $\zeta(0;\mathrm{Maxwell};S^4) = -31/45$.}
In Lorenz gauge with Faddeev--Popov ghosts the Maxwell Laplacian on transverse one-forms of the
unit $S^4$ has spectrum $\lambda_\ell = \ell(\ell+3)+2$ ($\ell\ge1$) with degeneracy
$d_\ell = 2\ell(\ell+3)(2\ell+3)/3$; the complex ghost contributes a scalar Laplacian with
spectrum $\ell(\ell+3)$. The combined zeta function evaluates at $s=0$ to
$\zeta(0;\mathrm{Maxwell};S^4) = -31/45 = -4\ag$, reflecting $\int_{S^4} E_4 = 64\pi^2$ and
$\int_{S^4}\avg{T^\mu{}_\mu} = 4\ag$.

\paragraph{R5: $\zeta(0;\mathrm{Maxwell};S^3) = 1$.}
On $S^3$ the transverse Maxwell spectrum is $\mu_\ell = (\ell+1)^2$ with degeneracy
$2\ell(\ell+2) = 2(\ell+1)^2 - 2$. With $n = \ell+1$,
$\zeta(s;S^3) = 2\zeta_H(2s-2,2) - 2\zeta_H(2s,2)$, and the Hurwitz values
$\zeta_H(-2,2) = -1$, $\zeta_H(0,2) = -\tfrac32$ give
$\zeta(0;\mathrm{Maxwell};S^3) = 2(-1) - 2(-\tfrac32) = 1$. Scalars and Dirac spinors give $0$;
the value $1$ is the photon-specific signature on $S^3$ that supplies the cap restoring force.

\paragraph{R8: DtN determinant $\det{}'\Lambda_{S^3}^{\mathrm{Maxwell}} = 1/(2\pi^3)$.}
The gauge-fixed Maxwell DtN is the combination
$\Lambda^{\mathrm{Maxwell}} = \Lambda^{\text{1-form}} - 2\,\Lambda^{\text{scalar ghost}}$ (transverse
one-form minus two Faddeev--Popov scalar ghosts), \emph{not} a single one-form spectrum. Its two
constituents are: the transverse one-form DtN with $\mu_\ell = \ell+1$, degeneracy $2\ell(\ell+2)$
($\zeta(0)=1$, $\zeta'(0) = 2\zeta_H'(-2,2)+\log 2\pi$); and the scalar-ghost DtN with
$\mu_\ell = \ell(\ell+2)/(\ell+1)$, multiplicity $(\ell+1)^2$ ($\zeta(0)=-1$,
$\zeta'(0) = \zeta_H'(-2,2)-\log\pi$, the same operator as R2). Hence
$\zeta(0;\mathrm{Maxwell}) = 1 - 2(-1) = 3$ (R4) and
$\zeta'_\Lambda(0) = [2\zeta_H'(-2,2)+\log 2\pi] - 2[\zeta_H'(-2,2)-\log\pi] = \log(2\pi^3)$
(the Hurwitz term $\zeta_H'(-2,2)$ cancels exactly), i.e.\
$\det{}'\Lambda_{S^3}^{\mathrm{Maxwell}} = 1/(2\pi^3) \approx 0.01613$, confirmed numerically to
$80+$ digits. (Pairing the one-form spectrum $\mu_\ell=\ell+1$ alone with this determinant is
incorrect: the one-form alone gives $\det{}' \approx 0.169$, not $1/(2\pi^3)$.)

\paragraph{R9: $\eta$-invariant $\eta(\Lambda_{S^3}^{\text{transverse 1-form}}) = 1$.}
The transverse one-form DtN eigenvalues $\mu_\ell = \ell+1 > 0$ are all positive, so its
eta-invariant equals its zeta value $\zeta(0;\Lambda^{\text{1-form}}) = 1$; it is conformally
invariant. (This is the one-form constituent of R8, $\zeta(0)=1$; the full gauge-fixed Maxwell
DtN has $\zeta(0)=3$, R4.)

\paragraph{R12: $\int_{D^4} Q_4 / \int_{D^4} E_4 = 1/4$.}
By Chang--Yang $\int_{D^4} Q_4 = 8\pi^2$ and by Gauss--Bonnet--Chern with boundary corrections
$\int_{D^4} E_4 = 32\pi^2$, giving the ratio $1/4$. The structural identity
$Q_4 = E_4/4 + (\text{conformal divergence})$ on conformally flat manifolds underlies the
resonance $4Q_4 = \mu_1(P_4) = 24$ on $S^4$.

\section{BFK decomposition and the DtN calculation}
\label{sec:bfk}
The Burghelea--Friedlander--Kappeler formula reads
$\log\det{}'\Delta_M = \log\det{}'\Delta_{M_+}^{\mathrm{Dir}} + \log\det{}'\Delta_{M_-}^{\mathrm{Dir}}
 - \log\det{}'\Lambda_\Sigma$. For $S^4 = D^4_+ \cup_{S^3} D^4_-$ with $D^4_+ \cong D^4_-$,
separation of variables in $S^3$ harmonics gives, for a transverse one-form harmonic with
$\Delta_{S^3}\omega_\ell = (\ell+1)^2\omega_\ell$, the radial problem
$[-\partial_\theta^2 - 2\cot\theta\,\partial_\theta + (\ell+1)^2/\sin^2\theta]u_\ell = \lambda u_\ell$.
The harmonic extension $u_\ell \propto (\sin\theta)^{\ell+1}$ yields the transverse one-form DtN
spectrum $\Lambda_{S^3}\omega_\ell = (\ell+1)\omega_\ell$ ($\zeta(0)=1$). The gauge-fixed Maxwell
determinant is then the combination $\Lambda^{\text{1-form}} - 2\Lambda^{\text{scalar ghost}}$,
giving $\zeta'_\Lambda(0) = \log(2\pi^3)$ and $\det{}'\Lambda_{S^3}^{\mathrm{Maxwell}} = 1/(2\pi^3)$
(R8); the one-form spectrum alone does \emph{not} give this determinant. This is consistent with the
$\zeta(0)$ bookkeeping: the scalar-ghost split uses the scalar DtN $\zeta(0)=-1$ and the
hemisphere value to give $2\zeta(0;D^4) - \zeta(0;S^4) = 2(-38/45) - (-31/45) = -1$
(bookkeeping (a) of Section~\ref{sec:spectral}), while the gauge-fixed Maxwell DtN itself has
$\zeta(0)=3$ (R4, bookkeeping (b)) --- two distinct operators that must not be conflated.

\section{The magnitude problem: Kaloper--Padilla sequestering}
\label{sec:sequester}

The first face of the cosmological constant problem --- why the gravitating vacuum energy is
$120$ orders of magnitude below the naive estimate --- is addressed by vacuum-energy
sequestering (Axiom A2, \cite{KP2014,BD2015}). In the Kaloper--Padilla construction the
classical action is supplemented by a global constraint that promotes the cosmological counterterm
to a dynamical variable. Integrating out that variable enforces
\begin{equation}
\label{eq:sequester}
\Lambda_{\mathrm{eff}} = \frac{1}{4}\,\frac{\int \sqrt{g}\,\avg{T^\mu{}_\mu}\,d^4x}{\int \sqrt{g}\,d^4x},
\end{equation}
i.e. the effective cosmological constant is the spacetime average of the matter stress-tensor
trace. The radiative quartic vacuum energy, being a spacetime constant, drops out of the
\emph{difference} between the local source and its global average: it does not gravitate. This
is the standard sequestering resolution of the magnitude problem.

\paragraph{Sequestering is internal to CSG, not an external add-on.} The constraint
\eqref{eq:sequester} need not be imposed by hand: it is generated by the conformal structure of
CSG itself. In Weyl geometry the gravitational action $S=\int\sqrt{\tilde g}\,(\alpha\tilde C^2 +
\beta\tilde R^2)$ decomposes, on fixing the Weyl gauge ($\tilde R = R + 6\nabla_\mu\omega^\mu +
6\omega_\mu\omega^\mu$), into a Weyl-squared term plus a scalar: defining
$\varphi \equiv \sqrt{2\beta}\,(R + 6\nabla\!\cdot\!\omega + 6\omega^2)$ turns the $R^2$ term into
$\tfrac12\varphi^2$, and integrating out the Weyl gauge field $\omega_\mu$ supplies its kinetic
term, giving $S=\int\sqrt g\,[\alpha C^2 + \tfrac12(\partial\varphi)^2 - V(\varphi) +
\mathcal{L}_{\mathrm{m}}]$. The compensator $\varphi$ is the Stückelberg field of the broken Weyl
symmetry, and it plays exactly the role of the Kaloper--Padilla dilaton: the Stückelberg
identity by which $\omega_\mu$ absorbs the scalaron is, term by term, the KP global constraint,
both being realizations of the same broken-CFT structure. Spontaneous breaking by the scalaron
VEV generates the Planck mass, $\Mpl^2 = v_\varphi^2/8\pi$, and fixes the residual
$\Lambda_{\mathrm{eff}} = V(v_\varphi)$; every radiative correction $\delta\Lambda$ is quadratic
in $\delta v_\varphi$ and reabsorbed by the scalaron at the same order in $\hbar$ (the linear
term is forbidden by the minimum condition), so $\Lambda_{\mathrm{eff}}$ is radiatively stable.
Sequestering is therefore \emph{derived within CSG} (the Built-in Sequestering proposition --- a
proof sketch by direct identification, with full all-orders rigor of the scalaron--dilaton map
and of the St\"uckelberg--KP correspondence in non-trivial backgrounds the remaining refinement),
not an adopted external mechanism; and because CSG carries no propagating graviton loop, it
evades the standard graviton-loop objection to KP. Crucially for what follows, Eq.~\eqref{eq:sequester} is also the seed of the value
problem: it expresses $\Lambda_{\mathrm{eff}}$ in terms of a global average that itself depends
on $\Lambda_{\mathrm{eff}}$ through the geometry, so the equation is implicit and must be solved
self-consistently. That self-consistency is the subject of the next section.

On the Euclidean hemisphere the same averaging, applied to the photon trace anomaly, yields the
numerator of the topological quotient: with $\int_{D^4} = \tfrac12\int_{S^4}$ and
$\int_{S^4}\avg{T^\mu{}_\mu} = 4\ag$ (Section~\ref{sec:anomaly}),
$\tfrac12\Lambda_{\mathrm{eff}}V_4(D^4) = \ag$, which combines with the $Q$-charge denominator
$\CQ = 8\pi^2$ in Section~\ref{sec:main}.

\section{Causal legitimation of the finite region and the saddle scale}
\label{sec:causal}

Two inputs above were stated as model-building choices: the restriction of the averaging to the
finite hemisphere $D^4$ rather than a global region requiring an observer measure, and the
identification of the saddle scale with $L_{\mathrm{HH}} = 1/H_0$. Both follow from a causal
variational principle. Shaw and Barrow \cite{ShawBarrow2011} promote $\Lambda$ to a field and
impose stationarity of the classical action under variations confined to the causal past of a
measurement --- the past light cone $M$:
\begin{equation}
\label{eq:SB}
\frac{dI_{\mathrm{class}}(\Lambda; M)}{d\Lambda} = 0 .
\end{equation}
The region $M$ needs no anthropic weighting: it is what a local measurement can access. The past
light cone of an observer in an open FLRW foliation is topologically a four-ball, so $M$ is the
Lorentzian counterpart of the Euclidean hemisphere $D^4$ on which the BFK decomposition operates.
The hemisphere is therefore not an arbitrary cutoff but the causal region itself, and the
boundary that BFK requires is the boundary of that region.

\paragraph{The saddle scale, parameter-free.} With extrinsic curvature $\mathrm{tr}\,K \sim H$ on
the bounding surface and four-volume $V_M \sim t_U\,A_{\partial M}$, Eq.~\eqref{eq:SB} fixes the
order of magnitude
\begin{equation}
\Lambda \sim \frac{H_0\,A_{\partial M}}{V_M} \sim t_U^{-2},
\end{equation}
so $L_{\mathrm{HH}} \sim 1/H_0$ with no free scale. Solving the self-consistency
$\Lambda = t_U^{-2}$ numerically gives $\OL \approx 0.49$ at factor-unity normalization, and at
the observed $\OL = 0.685$ the time ratio $t_\Lambda/t_U = 0.73$ --- the geometric origin of the
coincidence. The mechanism fixes the magnitude and the coincidence to order unity; it does
\emph{not} fix the precise value of $\OL$ (Section~\ref{sec:status}).

\paragraph{Structural identity with the cap mechanism.} Equation~\eqref{eq:SB} and the
cap-variational principle of Section~\ref{sec:cap} are one variational statement in two
realizations: a bulk trace drive balanced against a boundary restoring force over a finite causal
region. Shaw--Barrow is the Lorentzian, dynamical version, driven by the classical baryon trace
$\zeta_b = 1 - M_q/M_N \approx \tfrac12$; the cap mechanism is the Euclidean, topological version,
driven by the photon anomaly $3\ag$ against the DtN restoring force. The two differ decisively in
which trace enters: Shaw and Barrow set the photon contribution to the matter Lagrangian to zero
classically ($E^2 = B^2$ for radiation), which is exactly where the conformal anomaly --- the
quantum correction they omit --- resides. By the Komargodski--Schwimmer infrared selection
(Section~\ref{sec:anomaly}) the baryon trace decouples in the deep infrared, so the anomaly is the
surviving drive.

\paragraph{Why the realization must be topological.} A Lorentzian light-cone evaluation of the
anomaly fails quantitatively: the anomaly stress $\sim a^4 H^4$ is weighted to early times in the
light-cone integral, and the dynamical (Riegert/DtN-determinant) attempt to derive
$\ag/(8\pi^2)$ is too weak by the four-volume factor $(\Mpl/H)^4 \sim 10^{240}$
(Section~\ref{sec:uniqueness}, Mechanism~2). Only the topological evaluation --- the $a_4$
heat-kernel drive against the scale-invariant DtN $\zeta(0) = 3$, normalized by the Chang--Yang
charge $\CQ = 8\pi^2$ --- yields the parameter-free ratio. The causal principle legitimizes the
finite region and the scale; the topological character of the quotient is what makes the anomaly
evaluation finite and unique.

\section{The value problem: a self-consistency theorem}
\label{sec:theorem64}

Sequestering removes the radiative vacuum energy but leaves the residual scale to be fixed. We
now show that the residual $\Lambda$ is not free: it is the unique fixed point of the averaging
map~\eqref{eq:sequester}.

\subsection{The self-consistency iteration}
\begin{proposition}[Kaloper--Padilla self-consistency]
\label{thm:kp}
Let the effective cosmological constant be updated by the four-volume average of the total
energy density over a closed cosmological history,
$\Lambda^{(n+1)} = \tfrac14\avg{\rho_{\mathrm{total}}}^{(n)}_{V_4}$. Then on a closed cosmology
($k=+1$) with finite four-volume $V_4 = \int dt\, a^3 < \infty$, the update map is a contraction
with a unique fixed point $\Lambda^*$, and $\Lambda^*$ is the unique value satisfying the
trace-zero condition.
\end{proposition}
\begin{proof}
Two complementary parts. \emph{Convergence (analytic contraction).} Parametrize the update map
by $\Lambda$ via the Friedmann constraint at $a=1$, so that one iteration is $\mathcal{F}(\Lambda)
= \tfrac14\avg{\rho_{\mathrm{total}}(\Lambda)}_{\mathrm{ren}}$ with the four-volume ($a^3$-weighted)
average. Writing the trace function $G(\Lambda) \equiv \tfrac14\avg{\rho_{\mathrm{total}}}
- \Lambda = -\tfrac{\kappa}{4}\avg{T}$, the derivative is $\mathcal{F}'(\Lambda) = 1 + G'(\Lambda)$,
so the Lipschitz constant is $L = |1 + G'(\Lambda)|$. Increasing $\Lambda$ compresses the
closed history (shorter recollapse time, smaller $a_{\max}$), which strictly reduces the
$a^3$-weighted average of the positive matter and kinetic densities; hence $G$ is strictly
decreasing and $G'(\Lambda) < 0$. The contraction condition $L = |1 + G'(\Lambda)| < 1$ then holds
\emph{provided} $G'(\Lambda) > -2$. The sign $G'<0$ is established analytically (the compression
argument above); the bound $|G'| < 2$ is established \emph{numerically} ($|G'|\approx 0.7$, hence
$L\approx 0.13$--$0.18$ in the full pipeline, $L\approx 0.3$ in the reduced model). By the Banach
fixed-point theorem the iteration then has a unique attractor. A fully analytic bound on $|G'|$
remains the open refinement, so the contraction is \emph{numerically established} rather than
closed by an analytic proof. \emph{Correctness}: the
self-consistency condition $\Lambda^* = \tfrac14\avg{R}_{\mathrm{ren}}$ is, by the trace identity
$R = 4\Lambda + \kappa T$, algebraically equivalent to
\begin{equation}
\label{eq:tracezero}
\avg{T(\Lambda^*)}_{\mathrm{ren}} = 0 ,
\end{equation}
precisely the Kaloper--Padilla sequestering condition. Existence requires $k=+1$ with
$V_4 < \infty$; in flat or open cosmologies $V_4$ diverges and the volume average is
ill-defined. Since $\avg{-T}(\Lambda)$ is continuous, positive for small $\Lambda$ and negative
for large $\Lambda$ (the same monotonicity), the fixed point is unique by the intermediate value
theorem. The independent computation of Section~\ref{sec:m1} (M1-$\Lambda$) confirms the
trace-zero identity symbolically, the finiteness of $V_4$ for $k=+1$ against divergence for
$k=0,-1$, and the existence of a unique negative fixed point.
\end{proof}

\subsection{Compatibility of a closed history with open observation}
Proposition~\ref{thm:kp} requires a closed Euclidean structure with finite four-volume, which at
first sight conflicts with the open ($\Ok > 0$) Lorentzian patch predicted in
Section~\ref{sec:main}. There is no conflict, because three distinct regimes are involved:
(i) the \emph{full cosmological history} (Big Bang $\to$ expansion $\to$ recollapse $\to$ Big
Crunch), closed ($k=+1$), supplying the finite $V_4$; (ii) the \emph{present observable patch},
the Wick-rotated cap-shifted instanton, hyperbolic ($K<0$, $\Ok>0$) at the present epoch;
(iii) the \emph{Euclidean instanton} $D^4\subset S^4$, compact, on which the ratio
$\Ok/\OL = \ag/(8\pi^2)$ is fixed. These live on different manifolds related by analytic
continuation and fix complementary quantities; their coexistence is the content of the
no-boundary construction (Axiom A1).

\subsection{The Coleman--Weinberg vacuum and the status of the absolute value}
The fixed point can equivalently be read off the Coleman--Weinberg effective potential of the
scalaron, $V(\varphi) = A\,\varphi^4\big(\ln(\varphi^2/v_S^2) - \tfrac{25}{6}\big)$ with
$A \propto \ag$. Its unique nontrivial stationary point sits at $\varphi_{\min} = v_S\,e^{11/6}$
(verified in Section~\ref{sec:m1}, M1-CW), and in the stable convention ($V_{\min}<0$) this is
a minimum with negative $\Lambda^*$. What Proposition~\ref{thm:kp} establishes \emph{structurally}
is therefore that $\Lambda$ is not a free parameter: its sign is negative, the closed-history
requirement and the trace-zero characterization are proven, and the iteration converges.

What it does \emph{not} do unconditionally is pin the \emph{sharp} absolute value, and the
reason is instructive. The average is taken over $\rho_{\mathrm{total}}$, not $\rho_m$: the
scalaron potential component does not dilute, so $\avg{\rho_{\mathrm{total}}}_{V_4}$ is
\emph{$a_{\max}$-stable} (it tends to the non-diluting plateau as $a_{\max}\to\infty$), whereas
averaging $\rho_m$ alone would spuriously collapse as $1/a_{\max}^3$ --- a category error to be
avoided. This $a_{\max}$-stability of the total average and the topological constraint
$\Lambda_{\mathrm{eff}}V_4 = 4\ag$ are solid.

\paragraph{Correction: the volume route does not deliver the sharp value.} An earlier
sharpness analysis wrote $\Lambda^* = \tfrac14\avg{R}_{V_4}$ for a closed
matter+$\Lambda$ history and reported that a \emph{moderate} $a_{\max}$ breaks the de~Sitter
identity (giving $\avg{R}/4 \approx 0.76$--$0.89\times 3\OL$) while a long plateau returns the
identity $\Lambda_{\mathrm{eff}}=\Lambda$. That construction is \emph{geometrically
inconsistent} and its numbers are artefacts: it fixes a constant positive $\OL$ and then adds a
curvature term forcing $H^2=0$ at $a_{\max}$, but a universe with constant $\OL>0$ never
recollapses, so $H^2$ becomes \emph{negative} on $1<a<a_{\max}$ and was silently floored. The
correct treatment requires the genuine recollapse driver $V(\varphi)<0$. Integrating the full
Friedmann+Klein--Gordon system with the Coleman--Weinberg potential and feeding $-\Lambda^*/3$
back into the fixed point yields a sharper, and more sobering, picture:
\begin{itemize}
  \item $\Lambda^*$ is \emph{universal}: independent of $\varphi_0$ and of $\Omega_m$ (matter
  dilutes out of the $V_4$-weighted average). This universality is a genuine structural feature.
  \item Its \emph{value} is amplitude-set: $\Lambda^*\approx-0.10$ in $3H_0^2$ units for the
  fixed-point amplitude $A_{\mathrm{CW}}=\ag/128\pi^2$ (and $\approx-0.002$ for the M5 amplitude),
  \emph{not} the $-0.691$ quoted from the proxy.
  \item $\OL = (V(\varphi_0)-\Lambda^*)/3$ is \emph{not} fixed: it varies fully with $\varphi_0$.
  KP self-consistency pins $\Lambda^*$ but leaves $\varphi_0$ --- the scalaron value ``today'' ---
  free, so the absolute $\OL$ is not determined by this route.
\end{itemize}
The honest status is therefore: $\Lambda$ is not a free parameter and its sign is fixed
(structural, proven); the bulk vacuum cancels by sequestering, removing the $10^{122}$ hierarchy;
but the volume/self-consistency route delivers the universality \emph{structure} and the order of
magnitude only, neither the sharp $\Lambda^*$ nor the absolute $\OL$. Sharpening to
$\OL\approx0.704$ requires the halo-weighted prior or the causal closure of
Section~\ref{sec:measure}. None of this touches the theorem-level, prior-free output --- the
ratio $\Omega_K/\OL = \ag/8\pi^2$ --- which is independent of the absolute scale.

\subsection{The magnitude of the curvature: coherent spectral accumulation}
\label{sec:coherence}
The previous subsection settled what the volume/self-consistency route does and does not fix.
One question it deliberately left in words is the \emph{magnitude} question for the testable
output: after sequestering removes the bare $\Mpl^4$ vacuum, the renormalized photon trace
anomaly is still only a \emph{local} de~Sitter density
\begin{equation}
\rho_{\mathrm{anom}} = \tfrac14\avg{T^\mu{}_\mu} = \frac{\ag}{8\pi^2}\,3H^4 ,
\qquad
\frac{\rho_{\mathrm{anom}}}{\rho_{\mathrm{crit}}} = \frac{\ag}{8\pi^2}\frac{H^2}{\Mpl^2}
\sim 10^{-122},
\label{eq:localanom}
\end{equation}
so a naive reading would predict $\Ok \sim 10^{-125}$, not $\sim 10^{-3}$. A factor $(\Mpl/H)^2$
must lift the local density to the curvature scale. We record here the mechanism that supplies
it and state precisely how far it reaches.

\paragraph{Coherence is the linearity of the spectral sum.} The gauge-fixed Maxwell operator on
$dS_4$ has the discrete spectrum (Allen 1986; Higuchi 1991) $\lambda_n = n(n+2)H^2$ with
multiplicity $d_n = 2n(n+2)$, $n\ge1$. The vacuum sum runs over these \emph{fixed} eigenmodes,
$\sum_n d_n$, and is therefore \emph{linear} and deterministic: this linearity \emph{is} the
$\times N$ coherence. The familiar $\times\sqrt N$ estimate applies only to \emph{stochastic},
phase-uncorrelated contributions; the spectral modes on the compact sphere are phase-locked by
the geometry, so the $\zeta$-regularized sum over them is inherently coherent. The coherence is
thus not an added postulate but a property of the spectrum. The bulk mode count to the cutoff
$n_{\max}=\Mpl/H$ is
\begin{equation}
N_{\mathrm{tot}}(n_{\max}) = \sum_{n=1}^{n_{\max}} 2n(n+2)
= \tfrac23 n_{\max}^3 + 3n_{\max}^2 + \tfrac73 n_{\max} \sim \tfrac23\Big(\frac{\Mpl}{H}\Big)^3 ,
\end{equation}
and the horizon (area) projection reduces the volume scaling to
$N_{\mathrm{eff}} \sim (\Mpl/H)^2 = S_{dS}/\pi$, the de~Sitter entropy. The $(\Mpl/H)^2$
\emph{scaling} is rigorous; its $O(1)$ prefactor is convention-dependent (mode basis, projection
scheme).

\paragraph{What the coherent lift delivers.} Multiplying the local density~\eqref{eq:localanom}
by the coherent mode number gives, exactly,
\begin{equation}
\frac{N_{\mathrm{eff}}\,\rho_{\mathrm{anom}}}{\rho_{\mathrm{crit}}} = \frac{\ag}{8\pi^2},
\label{eq:coherentlift}
\end{equation}
i.e.\ precisely the dimensionless ratio $\Ok/\OL$. The $10^{122}$ gap is closed
\emph{structurally} for the quantity that the theory actually predicts and that DESI actually
measures: the coherent spectral sum lifts the local anomaly to the $O(10^{-3})$ curvature scale.
This is the magnitude resolution for the testable prediction, and it is independent of the
absolute-value question.

\paragraph{From the ratio to the absolute value.} Promoting the ratio~\eqref{eq:coherentlift} to
the absolute $\OL = 4\ag$ requires one further topological factor, and that factor is not free:
\begin{equation}
\frac{\rho_\Lambda}{\rho_{\mathrm{anom}}} = 32\pi^2\,\Big(\frac{\Mpl}{H}\Big)^2
= \Big(\int_{D^4} E_4\Big)\,N_{\mathrm{eff}},
\end{equation}
with $\int_{D^4}E_4 = 32\pi^2$ the hemisphere Euler number already used in the budget ladder.
Equivalently $\OL = |\zeta(0;\text{Maxwell};S^4)| = 4\ag$, the cutoff-free integrated trace
anomaly.

\paragraph{A second, independent anchor of the value: the conformal scale-flow.} The number
$4\ag$ is not only the spectral determinant $|\zeta(0)|$; it is also the conformal scale-anomaly
coefficient of the free energy. Under a global Weyl rescaling $g\to e^{2\sigma}g$ the
anomaly-induced (Riegert--Mottola) action obeys
\begin{equation}
\frac{\mathrm{d}S_{\mathrm{anom}}}{\mathrm{d}\sigma}
= \frac{\ag}{16\pi^2}\int_{S^4}\sqrt{g}\,E_4 = 4\ag ,
\end{equation}
because $\int_{S^4}E_4 = 64\pi^2$ is a topological invariant. Thus $4\ag$ is fixed
\emph{cutoff-free by two independent routes} --- the spectral determinant and the scale-flow ---
neither of which uses the $N_{\mathrm{eff}}$ bulk projection. Consequently the $O(1)$ prefactor
of that projection (the heuristic ambiguity $4/3\pi$ vs $1$) is \emph{irrelevant to the value}:
it is bypassed, not tuned. What remains hypothesis-level is only the \emph{identification} of
$\OL$ with this invariant. The scale-flow physically motivates it, but does not force it
dynamically: $S_{\mathrm{anom}}(\sigma)=4\ag\,\sigma$ is linear, so it has no $\sigma$-saddle, and
the dynamical fixing of $\sigma$ (the M5 route) returns $1/12\ag\approx0.48$, not $0.69$. The
identification is therefore the same conjecture flagged for A5 and in the budget ladder, now with
a physical anchor for its magnitude.

\paragraph{Honest status of this subsection.} \emph{Solved (structural):} the coherence (as the
linearity of the fixed-mode spectral sum), the $(\Mpl/H)^2$ scaling, and the magnitude of the
curvature ratio~\eqref{eq:coherentlift} --- the $10^{122}$ gap is closed for the testable $\Ok$
prediction; the bridge factor $32\pi^2=\int_{D^4}E_4$ is geometric, not a knob; and the
\emph{value} $4\ag$ is fixed cutoff-free by the two independent anchors above. \emph{Open (hypothesis-level, and not
decidable by DESI):} (a) the \emph{identification} $\OL=4\ag$ --- physically anchored by the
scale-flow but not forced by a saddle (the dynamical route gives $0.48$); and (b) the epoch
binding, since $\OL(z)$ is epoch-dependent, $4\ag$ is matched only at $z\approx0$, and no
dynamical attractor sits at $4\ag$ --- the coincidence problem~(iii), here inherited and
addressed only by observer-selection ($\OL\approx0.704$, Section~\ref{sec:measure}), not by a
timeless mechanism. The ratio $\Ok/\OL = \ag/8\pi^2$ remains untouched by both, as it is
independent of the absolute scale.

\section{The parameter-free relation: the eight-step derivation}
\label{sec:main}

We now assemble the topological quotient linking the residual cosmological constant to the
spatial curvature. The first six steps and the two Euclidean sub-steps of the seventh are
theorem-level; the one remaining genuinely chosen sub-step is the data-anchored reference state A5c.

\paragraph{Step 1 (anomaly integral).} By Eq.~\eqref{eq:anomaly} and Gauss--Bonnet--Chern,
$\int_{S^4}\avg{T^\mu{}_\mu} = \tfrac{\ag}{16\pi^2}\int_{S^4} E_4 = 4\ag$.

\paragraph{Step 2 (sequestering on the hemisphere).} By Eq.~\eqref{eq:sequester} and
$\int_{D^4} = \tfrac12\int_{S^4}$, $\tfrac12\Lambda_{\mathrm{eff}}V_4(D^4) = \ag$.

\paragraph{Step 3 ($Q$-charge).} By Lemma~\ref{lem:qcharge}, $\CQ = 8\pi^2$.

\paragraph{Step 4 (Wess--Zumino consistency).} Under $g \to e^{2\sigma} g$ the anomaly action
obeys the Wess--Zumino consistency relation \cite{Riegert1984}, with linear response
coefficient $\ag\,\CQ/(16\pi^2) = \ag/2$.

\paragraph{Step 5 (hemisphere quotient).} The dimensionless quotient
\begin{equation}
\label{eq:quotient}
\Rtop \equiv \frac{\ag}{\CQ} = \frac{\ag}{8\pi^2} = \frac{31}{1440\pi^2}
\end{equation}
contains no scale, no Newton constant, no Planck mass.

\paragraph{Step 6 (photon selection).} By Axiom A3 and the $26$-mechanism exhaustion
(Section~\ref{sec:uniqueness}), the coefficient is the photon's $\ag = 31/180$. This step uses
the BFK consistency check of Section~\ref{sec:spectral} but, as emphasized in
Section~\ref{sec:cap}, the derivation is otherwise \emph{boundary-free}.

\paragraph{Step 7 (Janus-point matching).} Three logically distinct sub-components:
\begin{itemize}[leftmargin=*,nosep]
\item[(7a)] \emph{Cap-saddle construction} (Section~\ref{sec:cap}): a variational principle on
$D^4$ produces a saddle at $\dtheta^* = \ag$ to leading order. Euclidean, geometric. \emph{Proven.}
\item[(7b)] \emph{Wick rotation} (Bucher--Goldhaber--Turok open-universe template
\cite{BGT1995,HawkingTurok1998}): translates the Euclidean cap-saddle into a Lorentzian
Friedmann patch with curvature $\Ok^{\mathrm{conformal}}$. \emph{Proven.}
\item[(7c)] \emph{Identification A5}: relates $\Ok^{\mathrm{conformal}}$ to the Friedmann
observable $\Ok^{\mathrm{obs}}$. Its Euclidean half A5a is a \emph{theorem} and the Wick-rotation
kinematics A5b are \emph{derived} given B1; the residual is the data-anchored reference choice
A5c. \emph{Specification (the only genuinely chosen element).}
\end{itemize}
Under (7a)+(7b)+(7c) the dimensionless Euclidean ratio projects onto the present-epoch
observable, giving Eq.~\eqref{eq:central}.

\paragraph{Step 8 (Friedmann translation).} With $\Om + \OL + \Ok = 1$ and
$\OL = \Lb\,|\Ok|$, $\Lb = 8\pi^2/\ag = 458.4$,
\begin{equation}
\label{eq:OK}
|\Ok| = \frac{1 - \Om}{\Lb \pm 1}, \qquad \OL = \Lb\,|\Ok|,
\end{equation}
the $+$ sign holding for an open universe. With the ladder-consistent budget $\Om = 0.3096$
(Section~\ref{sec:ladder}), $|\Ok| \approx 1.50\times10^{-3}$ and $\OL = 4\ag = 0.6889$, matching
the robust Planck~2018 (TT,TE,EE$+$lowE$+$lensing$+$BAO) value $\OL = 1-\Om = 0.6889$
($\Om = 0.3111\pm0.0056$)~\cite{Planck2018} exactly to four decimals; the CSG trio
$(\Om,\OL,\Ok) = (0.3096,\,0.6889,\,+0.0015)$ is consistent with Planck$+$BAO including its
curvature constraint $\Ok = 0.001\pm0.002$.

\begin{postulate}[A5: Anomaly--Curvature Correspondence]
The conformal-anomaly-derived curvature $\Ok^{\mathrm{conformal}}$ obtained from the hemisphere
quotient~\eqref{eq:quotient} via the cap-saddle and Wick rotation equals the observed Friedmann
spatial-curvature parameter $\Ok^{\mathrm{obs}}$ at the present epoch.
\end{postulate}
\noindent
A5 is not monolithic. Following the decomposition adopted here, its
Euclidean half --- A5a, the topological identity $\Ok^{\mathrm{conformal}} = \ag/\CQ$ --- is a
\emph{theorem} (Cauchy functional equation for linearity, Chern--Weil basis change for the
$1/8\pi^2$, and the APS index theorem on the de~Rham complex), up to the bounded BFK boundary
refinement $|\epsilon_\partial| \le 1.5\%$. What remains a physical postulate is the
identification A5b/A5c: that this Euclidean quantity, carried through the Wick rotation (A5b,
derivable in saddle-point approximation given B1), equals the present-epoch Friedmann observable
under the canonical $\Lambda$-dominated reference (A5c). Thus everything in
Eq.~\eqref{eq:central} except the A5b/A5c identification is theorem-level or literature-supported.

\section{The cap-variational mechanism, stability, and the boundary refinement}
\label{sec:cap}

\subsection{The saddle equation}
The hemisphere ``breathes'' by an angular shift $\dtheta$ of the polar cap. The photon anomaly
supplies a bulk drive linear in $\dtheta$ at the equator (Robin shift $\alpha = \tan\ag \approx
\ag$, Spectral Result R10), and the DtN determinant supplies a restoring force. Balancing them
gives the leading saddle $\dtheta^*_{\mathrm{lin}} = \ag$, and the exact angular dependence gives
\begin{equation}
\label{eq:saddle}
\ag\,\cos^4\dtheta = \sin\dtheta,
\qquad\Longrightarrow\qquad
\dtheta^*_{\mathrm{exact}} = 0.16391 ,
\end{equation}
a $\sim 5\%$ downward shift (cubic truncation $\dtheta^*_{(3)} = \ag - 2\ag^3$).

\subsection{Stability}
The second variation around the cap-saddle is positive in every inhomogeneous sector, and this
is established \emph{analytically}, not merely by mode-by-mode shooting. The deformed cap
$D^4(\theta_0 = \pi/2 + \dtheta^*)$ with $\dtheta^* = 0.164$ is a geodesic ball on $S^4$ with
$\theta_0 < \pi$, hence a proper subdomain. The fluctuation operator is the Hodge--de~Rham
Laplacian (with the fixed $S^4$ curvature potential, the same on cap and hemisphere), so its
Dirichlet eigenvalues obey \emph{domain monotonicity}: $\lambda_{\mathrm{Dir}}(\theta_0)$ is
strictly decreasing in $\theta_0$, with limit the closed-$S^4$ eigenvalue as $\theta_0\to\pi$.
The closed-$S^4$ inhomogeneous spectra are positive in closed form --- scalars $\ell(\ell+3)\ge4$,
coexact $1$-forms and transverse-traceless tensors $(\ell+1)(\ell+2)\ge6$ and $\ge12$
respectively --- so for every inhomogeneous mode
\[
\lambda_{\mathrm{Dir}}\big(\tfrac\pi2+\dtheta^*\big) \;>\; \lambda_{\mathrm{closed}}(S^4) \;>\; 0 .
\]
The shooting computation confirms the bracketing exactly (e.g.\ scalar $\ell=1$: hemisphere
$10.0 \to$ cap $8.16 \to$ closed limit $4$; TT $\ell=2$: $20.0\to17.27\to12$), but the positivity
itself follows from the two theorems (domain monotonicity and the positive closed spectrum),
not from the numerics. The transverse $1$-form is distinguished from the scalar by the
Hodge--Weitzenb\"ock term $\Delta^{(1)}_{\mathrm{Hodge}} = \Delta_{\mathrm{Bochner}} + 3$ on
$S^4$. The only non-positive direction is the constant ($\ell=0$) conformal/lapse mode, of Morse
index $1$, handled by the Gibbons--Hawking--Perry contour rotation (Section~\ref{sec:found});
the saddle is thus genuine. Reconfirmed independently in Section~\ref{sec:m1} (M1-H).

\subsection{The boundary refinement and why it cannot move the main result}
\label{sec:boundary}
The cap-shifted boundary carries Branson--Gilkey heat-kernel coefficients. Two facts make their
effect controlled. First, a \emph{rigorous} BFK symmetry: writing
$\zeta_{\mathrm{bnd}}(\pi/2+\delta) - 3/2 \equiv g(\delta)$, the gluing identity forces
$\zeta_{\mathrm{bnd}}(\pi/2+\delta)+\zeta_{\mathrm{bnd}}(\pi/2-\delta) = 3$ for all $\delta$, so
$g$ is an \emph{odd} function and all even-order Taylor coefficients vanish identically. Second,
and decisively, the main derivation (Steps 1--8) uses only the closed-$S^4$ anomaly integral
$\ag$ and the $Q$-charge $\CQ = 8\pi^2$; neither couples to boundary heat-kernel coefficients.
The boundary appears \emph{only} in the BFK consistency check (Step 6) and in the cap-saddle
shift. Consequently the worst-case boundary uncertainty is geometrically excluded from the
prediction chain $\ag \to \Ok$. Quantitatively, in the natural renormalization scheme
($\mu = R^{-1}$ fixed) the shift is $0\%$ (BFK-protected); in the cap-scale scheme it is bounded
at $|\Delta\Ok|/\Ok \le 1.5\%$; the unprotected worst case ($\le 49\%$) does not propagate,
being excluded by the boundary-free structure of the main derivation.

\subsection{The prediction hierarchy L0--L3}
\begin{table}[h]
\centering
\renewcommand{\arraystretch}{1.3}
\begin{tabular}{@{}c l l l@{}}
\toprule
\textbf{Level} & \textbf{Formula} & \textbf{Value} & \textbf{Status} \\
\midrule
L0 & $R_0 = \ag/(8\pi^2)$ & $2.181\times 10^{-3}$ & topological-anomalic (exact theorem) \\
L1 & $R_{\mathrm{cap}} = \dtheta^*_{\mathrm{exact}}/(8\pi^2)$ & $2.076\times 10^{-3}$ & cap-variational refined \\
L2 & $R_{\mathrm{cap}+\partial} = R_{\mathrm{cap}}(1+\epsilon_\partial)$ & $2.05$--$2.11\times 10^{-3}$ & boundary check ($|\epsilon_\partial| \le 1.5\%$) \\
L3 & $\Ok^{\mathrm{obs}} = R\cdot\OL$ & $1.42$--$1.49\times 10^{-3}$ & present-epoch observable \\
\bottomrule
\end{tabular}
\caption{The four-level prediction hierarchy. L0 is the exact topological identity on $D^4$. L1
incorporates the exact cap-saddle solution of Eq.~\eqref{eq:saddle}. L2 is the boundary
consistency band (Section~\ref{sec:boundary}); it does \emph{not} enter the boundary-free L0/L1
derivation. L3 is the Friedmann translation.}
\label{tab:hierarchy}
\end{table}
\noindent
Headline statements are quoted at L0; numerical data comparison uses L1 or L3. The L0-vs-L1
split, $\Delta R = 1.05\times 10^{-4}$ ($\Delta\Ok = 0.72\times 10^{-4}$), is a small refinement
whose observational discrimination is forecast in Section~\ref{sec:obs}.

\section{Technical refinements and foundational inputs}
\label{sec:refine}

\paragraph{Boundary heat-kernel coefficients (severity low).}
The Branson--Gilkey--Kirsten--Vassilevich boundary coefficients for the cap-shifted $S^3$ enter
\emph{only} the consistency checks, not the main derivation, which lives on the closed $S^4$. The
BFK gluing identity $\zeta_M(0;D^4(\theta_0)) + \zeta_M(0;D^4(\pi-\theta_0)) = 104/45$ holds for
all $\theta_0$; decomposing $\zeta_M = \zeta_{\mathrm{bulk}} + \zeta_{\mathrm{bnd}}$ and writing
$\delta = \theta_0 - \pi/2$, the bulk part satisfies
$\zeta_{\mathrm{bulk}}(\pi/2+\delta) + \zeta_{\mathrm{bulk}}(\pi/2-\delta) = -4\ag$ and the
boundary part $\zeta_{\mathrm{bnd}}(\pi/2+\delta) + \zeta_{\mathrm{bnd}}(\pi/2-\delta) = 3$
(both constant in $\delta$). Hence $g(\delta) := \zeta_{\mathrm{bnd}}(\pi/2+\delta) - 3/2$ is an
\emph{odd} function: all even-order Taylor coefficients vanish identically. This is a rigorous
theorem, not a perturbative estimate. The leading bulk shift is $|\Delta\zeta_{\mathrm{bulk}}|
= 3\ag^2 \approx 0.088$; in the BFK-protected natural scheme the boundary shift is zero, and in
the cap-scale scheme it is bounded by $|\Delta\Ok|/\Ok \lesssim 1.5\%$ (the L2 band of
Table~\ref{tab:hierarchy}).

\paragraph{Scheme independence.}
The leading prediction $\Rtop = \ag/(8\pi^2)$ is built from the trace-anomaly integral (a
renormalization-group invariant) and the $Q$-charge (a topological invariant); neither depends on
the renormalization scale $\mu$, so $\Rtop$ is scheme-independent. The natural scheme
($\mu = R^{-1}$) is adopted throughout; the cap-scale scheme shifts only subleading corrections,
by $\lesssim 1.5\%$.

\paragraph{Picard--Lefschetz contour analysis.}
The saddle equation $\ag\cos^4 z = \sin z$ has one real saddle ($\dtheta^* = 0.16391$,
$\mathrm{Im}\,S = 0$ exactly); every complex saddle in the box $|\mathrm{Re}\,z| < 8$,
$|\mathrm{Im}\,z| < 2$ (sixteen of them) satisfies $|\mathrm{Im}\,S| \ge 0.78$. By
Picard--Lefschetz theory no Stokes line connects the real saddle to any complex saddle --- the
gap exceeds the connection threshold by a factor $>7\times 10^5$ --- so the steepest-descent
thimble through the cap is topologically isolated and $n_{\mathrm{complex}} = 0$ is established
\emph{analytically} for the reduced cap-action, not merely numerically; the RK45 thimble flow
confirms it independently. The cap saddle therefore dominates the reduced cap-action contour
integral. The embedding of this result into the full infinite-dimensional no-boundary path
integral is assembled semiclassically (Section~\ref{sec:found}): the cap-plus-scalar enumeration
is complete, and the prediction $\ag/(8\pi^2)$ is all-orders stable as a ratio of
$\hbar$-independent invariants. As a conditional fallback we retain the worst-case estimate ---
were a complex saddle nonetheless to contribute ($|n_{\mathrm{complex}}| = 1$ in the full theory),
the impact would be $|\epsilon_{\mathrm{PL}}| \le 6.6\times10^{-4}$ in $\Ok^{\mathrm{obs}}$ ---
about $44\%$ of the prediction, but below DESI DR2 sensitivity and within DESI DR3 / Euclid reach.

\paragraph{Foundational computations on the three residual points.}
Three explicit calculations sharpen the residual open items. \emph{(i) The embedding $B1'$.} The
Euclidean fluctuation Hessian was computed by shooting in all three transverse sectors (scalar,
co-exact $1$-form, transverse-traceless tensor) for $\ell \le 12$ on the deformed cap; every
eigenvalue is strictly positive (global minimum $8.16$), and the centrifugal term
$\ell(\ell{+}k)/\sin^2\theta$ grows monotonically, so positivity for all higher $\ell$ is
guaranteed asymptotically. The cap saddle is therefore a Morse-index-$1$ critical point (the
single negative direction is the lapse), enlarging the transverse mode space adds only convergent
Gaussian factors, and $n_{\mathrm{complex}} = 0$ is preserved at one loop for any finite
truncation. The genuine residual is non-perturbative (infinite-DoF, non-Gaussian), suppressed by
$A_{\mathrm{M5}} \approx 1.2\times10^{-6}$ --- a limitation generic to all of semiclassical
quantum gravity, not specific to this framework. \emph{(ii) The unconditional sign.} The
minisuperspace lapse on-shell action was derived symbolically for two boundary data: Dirichlet
$q(0)=0$ yields four saddles (including a complex pair with opposite-sign weights, the
Feldbrugge--Lehners--Turok ambiguity), whereas Neumann $\dot q(0)=0$ yields two real saddles --- the
analytic structure itself is boundary-condition-dependent, and the anomaly term
($\ag/2 \approx 0.086$ against $6$) does not lift it. The sign is thus fixed only once a
no-boundary boundary condition (equivalently, an integration contour) is chosen; this is a
foundational definitional input of quantum cosmology, not a CSG-KP-specific gap. \emph{(iii) The
reference epoch A5c.} The physical density ratio obeys $\Ok/\OL = (\Ok^0/\OL^0)(1+z)^2$ (since
$\rho_K \propto a^{-2}$, $\rho_\Lambda = \text{const}$), ranging from $2.18\times10^{-3}$ today
to $\sim 2.6\times10^3$ at recombination. The prediction $\ag/(8\pi^2)$ is specifically the $z=0$
value; a reference epoch is therefore genuinely required and data-anchored, confirming that A5c
is irreducible. The redshift-constancy established earlier is a property of the Euclidean
conformal quantity, not of the physical density ratio.

\paragraph{The fixed-point value: reconciliation.}
Two legacy pipelines reported $\Lambda = -0.81331$ and $\Lambda^* \approx -0.691$. Both are
\emph{superseded} by the correction of Section~\ref{sec:theorem64}: the $\langle R\rangle/4$
sharpness proxy on which they rested is a floored-geometry artefact (a constant $\OL>0$ does not
recollapse, so the forced turnaround drove $H^2<0$ and was silently floored). The corrected
Friedmann$+$Klein--Gordon dynamics give a \emph{universal} fixed point $\Lambda^*$ (independent of
$\varphi_0$ and $\Omega_m$) but only \emph{amplitude-set}: $\Lambda^*\approx-0.10$ in $3H_0^2$ units
for $A_{\mathrm{CW}}=\ag/128\pi^2$ ($\approx-0.002$ for the M5 amplitude), \emph{not} the legacy
$-0.691$. What survives is structural and unchanged: $\Lambda$ is not free, its sign is negative,
the bulk vacuum cancels, and the iteration is a Banach contraction (constant $L\approx0.18$--$0.30$,
sign $G'<0$ analytic, $|G'|<2$ numerical). What does \emph{not} survive is a sharp absolute value:
$\OL=(V(\varphi_0)-\Lambda^*)/3$ leaves $\varphi_0$ free, so the homogeneous route fixes neither the
sharp $\Lambda^*$ nor the absolute $\OL$. None of this bears on the ratio $\Ok/\OL=\ag/8\pi^2$.

\paragraph{Conditional predictions.}
If the lightest neutrino is exactly massless, the Weyl anomaly $a_W = 11/720$ adds to the photon,
giving $|\Ok| = 0.00162$ ($\Delta|\Ok| = 1.3\times10^{-4}$ from photon-only; Euclid+CMB-S4
discriminates at $\sim0.5\sigma$). A second massless $U(1)$ would double the coefficient to
$2\ag = 62/180$. Both are falsifiable by DESI DR3 / Euclid.

\paragraph{Foundational inputs.}
The smooth Janus point at the equator depends on Axiom A1 (Hartle--Hawking no-boundary); the
framework is compatible with alternative Janus-providing scenarios (Shape Dynamics, Boyle--Turok
CPT-bounce) but A1 is the most-studied choice. The coefficient $\ag = 31/180$ assumes the
Bunch--Davies vacuum, the unique maximally symmetric Hadamard state on de Sitter; non-trivial
alpha-vacua would modify it but are physically disfavoured without new structure.

\section{Uniqueness: the classification and the falsified mechanisms}
\label{sec:uniqueness}

\begin{proposition}[Uniqueness of the invariant]
\label{prop:unique}
Among dimensionless combinations of conformal-anomaly and $Q$-curvature data on $D^4$, the
ratio $\ag/\CQ = \ag/(8\pi^2)$ is the unique candidate satisfying: (i) scale invariance (no
$\Mpl$, no $G$); (ii) conformal naturalness ($Q_4$, not $E_4$, as denominator); (iii) photon-only
infrared content (Komargodski--Schwimmer selection); (iv) constancy of the action-parameter
ratio $3|K|/\Lambda$ (epoch-independent as a ratio of fixed FRW constants, the density parameters
redshifting as $1/a^2$).
\end{proposition}

\noindent
\emph{Rigorous core.} The decisive structural fact is exact: on $S^4$ the Paneitz operator $P_4$
has kernel equal to the constants, so the constant conformal mode $\sigma_0$ is its \emph{unique}
zero mode, and $\int_{D^4} Q_4$ is the unique linear functional appearing as its Wess--Zumino
source. Any alternative normalization either redefines $\sigma_0$ by a constant (leaving the ratio
invariant) or is not the conjugate quantity; the Chang--Yang relation $\int Q_4/\int E_4 = 1/4$
provides an independent cross-check ($4\ag/32\pi^2 = \ag/8\pi^2$). \emph{Honest limit.} Criteria
(i)--(iv) are each implied by an axiom and are \emph{sufficient} to select $8\pi^2$ from the
natural candidate list; we do not claim they are minimal --- a systematic proof that no broader
admissible criterion set selects a different invariant is open, and the uniqueness is therefore
established within the conjugate-pair plus (i)--(iv) framework, not as an unconditional theorem
over all conceivable constructions.

A complementary constructive argument exhausts the dynamical alternatives: every attempt to
obtain $\Ok^{\mathrm{conformal}} = \ag/(8\pi^2)$ from a \emph{local field equation} rather than
the topological quotient fails, organizing into six no-go categories --- (I) conformal
mechanisms linear in $\sigma_0$ or vanishing; (II) Lorentzian dynamics reintroducing Newton's
constant; (III) massive fields regenerating the $10^{120}$ hierarchy; (IV) squashing
deformations only quadratic; (V) spectral invariants conformally protected; (VI) vacuum energy
absorbed by sequestering. Table~\ref{tab:mechanisms} lists the twenty-six explicit mechanisms.

\begin{longtable}{@{}c p{6.4cm} c p{4.2cm}@{}}
\caption{The twenty-six falsified mechanisms for deriving
$\Ok^{\mathrm{conformal}} = \ag/(8\pi^2)$ from local dynamics.}\label{tab:mechanisms}\\
\toprule
\# & Mechanism & No-Go & Failure mode \\ \midrule \endfirsthead
\toprule \# & Mechanism & No-Go & Failure mode \\ \midrule \endhead
1 & Riegert variational & I & linear in $\sigma_0$; no $\sigma_0^2$ \\
2 & DtN determinant (naive) & I & too weak by $10^{240}$ \\
3 & KP stiffness & I & conformal sector unchanged \\
4 & Gauss--Bonnet (conformal) & I & identically zero on FRW \\
5 & Gauss--Bonnet (geometric) & I & cubic in $a$; no $\sigma_0^2$ \\
6 & Resonance $4Q_4 = \mu_1$ & I & coupling coefficient $\approx 0$ \\
7 & M\"obius nonlinearity & I & no 4D solution \\
8 & Topology sum & I & not regularizable \\
9 & Conformal Perelman flow & I & $P_4$ remains covariant \\
10 & Two-loop stiffness & I & zero for free conformal fields \\
11 & $L=1$ coupled system & I & $\sigma_0^*$ is $\ag$-independent \\
12 & Counter-field & I & symmetric hemispheres cancel \\
13 & $C_Q$-decoherence & II & requires $K = \CQ^2$ input \\
14 & $Q$-prescription & II & restoring force $=0$ after KP \\
15 & Topology path integral & I & zero-mode divergence \\
16 & Fokker--Planck for $\sigma_0$ & II & gives $\Ok \propto 1/\ag$ (wrong) \\
17 & Lorentzian field equation & II & $(H/\Mpl)^2$ suppression \\
18 & Anomalous field equation & IV & squashing only quadratic \\
19 & Conformal scalar VEV & III & $V(\sigma_0) = \text{const}$ \\
20 & Coleman--Weinberg & I & $\beta$-function $=$ anomaly \\
21 & Explicit mass $m^2$ & III & hierarchy of $10^{249}$ \\
22 & Topological mass $\xi R$ & III & wrong sign \\
23 & $\zeta'(0)$ spectral & V & transcendental constant \\
24 & Lichnerowicz spectral & V & no $\sigma_0$ dependence \\
25 & Heat-kernel spectral & V & rigid on $S^4$ \\
26 & Stochastic & VI & $\avg{C^2}$ absorbed by KP \\
\bottomrule
\end{longtable}
The lesson is methodological: the correct route is not a new field equation but the topological
hemisphere quotient, with the cap-saddle providing the geometric realization.

\section{Historical methodology}
\label{sec:history}
Five reference cases share one logical pattern, \emph{compact space $+$ local condition
$\Rightarrow$ global parameter}: Balmer$\to$Bohr 1913 (electron on a compact configuration space,
single-valuedness selects discrete levels); Dirac 1931 ($eg = n\hbar/2$ from $S^1$
single-valuedness around the string); the quantum Hall effect 1980 ($\sigma_{xy} = ne^2/h$ from
the integer Chern number on a torus); Hawking--Gibbons 1977 ($T = \kappa/2\pi$ from regularity of
the compact Euclidean section); and Casimir 1948 ($\rho = -\pi^2/720d^4$ from boundary conditions
on a compact transverse interval). CSG--KP has the compact space ($D^4$) and the conditions (KP
sequestering, photon-uniqueness, Janus matching) but, by the $26$-mechanism exhaustion
(Section~\ref{sec:uniqueness}), no \emph{local} mechanism that singles out a value: there is no
Hawking conical regularity (the equatorial $S^3$ is maximal, not conical), no Dirac quantization
(the conformal connection is trivial), no Hall index ($\eta$ is conformally invariant), no
Casimir extremum, and no anomaly-matching obstruction (the FLRW slice is conformally flat
throughout). The prediction therefore belongs to a distinct category --- a topological/anomalic
ratio, not a regularity-quantized observable --- which is the structural content of the
complementarity argument.

\section{The coincidence problem: structure-formation averaging}
\label{sec:coincidence}

The third face of the cosmological constant problem --- why $\OL$ and $\Om$ are comparable now
--- is not solved by the ratio~\eqref{eq:central}, whose constant object is the action-parameter
ratio $\mathcal{R}=3|K|/\Lambda$ (the density-parameter ratio itself redshifts as $\mathcal{R}/a^2$).
A clean
physical constraint narrows the solution space: a dynamical dark-energy component cannot
\emph{simultaneously} have $w \approx -1$ (data) and track matter (coincidence); the two are
mutually exclusive, and sequestering, being a constant subtraction, leaves $w_{\mathrm{eff}} =
-1$. The coincidence is therefore resolved not by dark-energy \emph{dynamics} but by the
\emph{averaging prescription} for $\Lambda$, following Lombriser's structure-formation
counterpart to Kaloper--Padilla \cite{Lombriser2019}.

Lombriser's mechanism replaces the homogeneous four-volume average of
Section~\ref{sec:theorem64} by an average weighted toward gravitationally bound (halo) regions,
evaluated with the spherical-collapse top-hat evolution
\begin{equation}
\label{eq:tophat}
y'' + \Big(2 + \frac{H'}{H}\Big) y' + \tfrac12 \Om(a)\,(y^{-3} - 1)\,y = 0,
\qquad ' \equiv \frac{d}{d\ln a},
\end{equation}
with $y = (\rho_m/\bar\rho_m)^{-1/3}$ and matter-dominated initial data $y_i = 1-\delta_i/3$,
$y_i' = -\delta_i/3$. The single additional assumption is a uniform prior on $y$, i.e.
$y(t_0) = \tfrac12$. With the equality normalization ($a_{\mathrm{eq}} \equiv 1$,
$\Om(a) = a^{-3}/(a^{-3}+1)$) the critical (longest-lived) shell has overdensity
$\delta_i^{\mathrm{crit}} = 1.134\times 10^{-3}$ and reaches $y = \tfrac12$ at
$a_0/a_{\mathrm{eq}} = 1.336$, giving
\begin{equation}
\label{eq:omlambda}
\OL = \frac{1}{a_0^{-3} + 1} = 0.7044 .
\end{equation}
This value is \emph{derived}, not calibrated --- but conditional on the uniform-prior
assumption $y(t_0) = \tfrac12$, which is anthropic-like and is not itself derived from the
framework. The homogeneous and halo-weighted averages are two faces of the same
$\Lambda = \tfrac14\avg{R}$ structure, differing in averaging prescription. The
\emph{homogeneous} (volume) route delivers, as corrected in
Section~\ref{sec:theorem64}, the universality \emph{structure} of $\Lambda^*$ and the order of
magnitude only --- not the absolute scale, since $\varphi_0$ remains free. The
\emph{halo-weighted} route returns $0.704$ but at the cost of the anthropic-like
$y(t_0)=\tfrac12$. The Euclidean topology fixes the dimensionless \emph{ratio} $\Ok/\OL$
rigorously; the absolute \emph{scale} is fixed by neither the volume route (free $\varphi_0$)
nor without a prior by the halo route. The resulting trio,
\begin{equation}
\OL = 0.7044, \qquad \Ok = \frac{\ag}{8\pi^2}\,\OL = 1.54\times 10^{-3}, \qquad \Om = 0.2956 ,
\end{equation}
is internally consistent and within $2.8\%$ ($\OL$) and $0.7\sigma$ ($\Ok$) of observation.

\subsection{A geometric budget ladder (conjecture)}
\label{sec:ladder}
A second, sharper route to the absolute scale is suggested by the numbers themselves, and is
recorded here \emph{as a conjecture}, with its open step made explicit. The factors that appear
across the framework form a $2^n\pi^2$ ladder of $4$D conformal geometry: the one-loop factor
$16\pi^2 = 2^4\pi^2$ in $\avg{T} = \ag E_4/16\pi^2$; the hemisphere Euler integral
$\int_{D^4}E_4 = 32\pi^2 = 2^5\pi^2$ ($\chi=1$); and the full-sphere integral
$\int_{S^4}E_4 = 64\pi^2 = 2^6\pi^2$ ($\chi=2$). Within this ladder the three budget quantities
all follow from $\ag$ and the Euler integrals:
\begin{align}
\OL &= \int_{S^4}\avg{T}\sqrt{g} = 4\ag = 0.6889 , \label{eq:ladder-ol}\\
\frac{\Ok}{\OL} &= \frac{4\ag}{\int_{D^4}E_4} = \frac{\ag}{8\pi^2} = 2.181\times10^{-3} ,\\
\Ok &= \frac{\OL^2}{\int_{D^4}E_4} = \frac{\ag^2}{2\pi^2} = 1.503\times10^{-3} .
\end{align}
The same numerator $4\ag$ appears in the leading term and in the ratio; dividing by
$\int_{D^4}E_4 = 32\pi^2$ turns the full-sphere anomaly into the cap ratio. This yields the exact
algebraic identity $\Ok\cdot 32\pi^2 = \OL^2$, and a \emph{perturbation hierarchy in $\ag$}:
dark energy is the leading anomaly effect ($\OL\sim\ag^1$) while curvature is second order
($\Ok\sim\ag^2$). The smallness $|\Ok|\ll\OL$ is then structural --- curvature is quadratically
suppressed in the anomaly --- not a tuning, mirroring the holographic $V_4 = A_2^2$ relation of
Section~\ref{sec:psf}. The split is geometric: the full $S^4$ (the timeless no-boundary
substrate) carries the absolute $\OL$, the hemisphere $D^4$ (the observer past light cone)
carries the ratio, and $32\pi^2 = \int_{D^4}E_4$ is the substrate-to-observer bridge. The full
budget $\OL = 0.6889$, $\Om = 0.3096$, $|\Ok| = 1.5\times10^{-3}$ matches Planck~2018 (base
$\Lambda$CDM) at $0.0\sigma$ ($\OL$), $-0.27\sigma$ ($\Om$), consistent ($\Ok$), from
$\ag = 31/180$ and $\chi=2$ alone.

\paragraph{Status (the one open step).} The ladder \emph{algebra} is exact and the factor
$32\pi^2$ is geometric, not a free normalization; the ratio $\ag/8\pi^2$ is theorem-level. What
is \emph{not} derived is the identification~\eqref{eq:ladder-ol} itself: why the dark-energy
density \emph{fraction} should equal the integrated anomaly \emph{action} on the full sphere. The
local energy-density route ($\rho_{\mathrm{anom}}/\rho_{\mathrm{crit}}\times N_{\mathrm{eff}}$, the
de~Sitter identity of Section~\ref{sec:psf}) yields only the \emph{ratio} $\ag/8\pi^2$, not
$4\ag$; an action/topological identification of the dark energy is required, which A5 does not
supply (A5 establishes the ratio only), and which is in tension with the KP sequestering of the
$\Lambda$-channel. The numerical match is therefore striking and the structure coherent, but
\eqref{eq:ladder-ol} is a conjecture pending this action-identification bridge --- not a proven
prediction of the absolute scale.

\paragraph{Refinement (this work): the prefactor is computed, the residue localized, the scope confined.}
Three results sharpen the open step without closing it. \emph{(i) The lift prefactor is computed.}
The factor identification~\eqref{eq:ladder-ol} requires, $\rho_\Lambda/\rho_{\mathrm{anom}} =
32\pi^2(\Mpl/H)^2$, is exactly the Bekenstein--Hawking horizon degree-of-freedom count $A/G = 4\,S_{dS}$
--- a standard holographic quantity, not a free normalization --- and with it $\OL =
(A/G)\,\rho_{\mathrm{anom}}/\rho_{\mathrm{crit}} = 4\ag$ falls out, the coherent ($\times N$)
accumulation being the unique scale-free choice (incoherent $\sqrt N$ gives an $H$-dependent value). \emph{(ii) The residual assumption is localized.}
An in-in (Schwinger--Keldysh) evaluation shows the lift is not an in-in mode sum (that is already
inside $\rho_{\mathrm{anom}}$) but the Cohen--Kaplan--Nelson UV--IR bound $\rho_\Lambda\le\Mpl^2/L^2$:
its Schwarzschild-consistency part \emph{follows} from gravity and fixes the $10^{-122}$ magnitude,
while its \emph{saturation} is the assumption P5 --- now a named holographic principle, not a vague
accumulation. \emph{(iii) The scope is confined.}
This residue sits in the absolute value alone: the ratio $\Ok/\OL = \ag/8\pi^2$ is scale-invariant
($H$ and $\Mpl$ cancel identically), so the falsifiable Euclid prediction is independent of the
unproven saturation; deriving the saturation from scale-invariance is a fallacy, since scale-invariance
constrains only the ratio. The conjecture is thereby
sharpened --- prefactor computed, principle named (CKN saturation), scope limited to the conjectural
bonus --- but not proven.

\paragraph{A no-go: the value is an action invariant, not a dynamical equilibrium.} The open step
above is structural, not a missing calculation. Every dynamical or energy-density route to the
absolute $\OL$ gives a \emph{different} number, and none gives $4\ag$:
\begin{itemize}[topsep=2pt,itemsep=1pt]
  \item KP sequestering $\rho_\Lambda = \int\avg{T}/V_4$ --- $(\ag/8\pi^2)(H/\Mpl)^2 \sim 10^{-122}$;
  \item the canonical Euclidean vacuum energy $\rho_\Lambda = -\,\mathrm{d}S_E/\mathrm{d}V_4$ on the
  $S^4$ saddle --- $\tfrac12$ (the gravitational term dominates; the anomaly enters at $10^{-122}$);
  \item the M5/KP fixed point $1/12\ag \approx 0.484$; the Coleman--Weinberg iteration ($\varphi_0$
  free); the Riegert mode (linear, no saddle);
  \item the causal closures over the forced apparent-horizon surface --- bare dimensional choices
  with median $\approx0.93$, KP stationarity $\avg{R}_M/4\approx14$, and the Lorentzian light-cone
  anomaly too weak by $(\Mpl/H)^4\sim10^{240}$.
\end{itemize}
The \emph{only} route to $4\ag$ is the pure action invariant $\int_{S^4}\avg{T}\sqrt{g} = 4\ag$
($=|\zeta(0)|=$ the Riegert scale-flow), cutoff-free and exact but an \emph{action}, not a
dynamically forced density. Hence $4\ag$ is not a dynamical equilibrium value at all; it is a
topological--spectral invariant, and the identification of the dark-energy fraction with it cannot
be derived from any dynamics --- each dynamics demonstrably yields something else. This is the
precise structural reason the step is open.

\paragraph{The sharpest form of the conjecture.} The identification $\OL = 4\ag$ is exactly
equivalent to a single normalization,
\begin{equation}
\label{eq:sdscq}
N_{\mathrm{eff}} = \frac{S_{dS}}{\CQ}
= \frac{8\pi^2\Mpl^2/H^2}{8\pi^2} = \Big(\frac{\Mpl}{H}\Big)^2 ,
\end{equation}
the Gibbons--Hawking de~Sitter entropy per $Q$-curvature charge. With this $c=1$ normalization the
coherent lift of Section~\ref{sec:coherence} returns $\OL = 4\,c\,\ag = 4\ag$. The open question
thereby narrows from ``a vague action identification'' to ``why this one geometrically natural
normalization (entropy per $Q$-charge)''. It is sharper and more falsifiable, but $S_{dS}/\CQ$ is
not more forced than the bulk-projection prefactor $4/3\pi$, so it remains a conjecture --- now in
its tightest form.

\subsection{Why the sharp value is not pinned by the history ensemble: the measure problem}
\label{sec:measure}
The conditional status of $\OL = 0.704$ deserves a sharper statement, because two independent
attempts to remove the uniform-prior assumption were carried through and both fail for the same
categorical reason. First, Lombriser and CSG--KP are the \emph{same constraint class}: Lombriser
varies the Einstein--Hilbert action with respect to the Planck mass, producing a global
Lagrange-multiplier constraint on the averaged Ricci scalar $\avg{R}$ --- structurally identical
to $\Lambda^* = \tfrac14\avg{R}_{\mathrm{ren}}$ of Proposition~\ref{thm:kp}. The single free element
is the prior on the observer's position in structure formation.

\paragraph{The collapse-time lever.} The value $0.704$ is the limit of \emph{infinitely
long-lived} shells ($a_{\mathrm{collapse}} = 6686$ at equality normalization); for a finite
maximum shell age $\OL$ decreases, and the observed $\OL = 0.685$ is reached already at
$a_{\mathrm{collapse}} \approx 8.67$. CSG--KP appears to supply exactly what Lombriser lacks:
Proposition~\ref{thm:kp} forbids $t_{\max}\to\infty$, since a finite four-volume requires a background
recollapse driven by the Coleman--Weinberg potential, which would truncate the shell lifetime and
pull $0.704$ toward observation. \emph{The lever fails quantitatively.} Forward integration of the
full background dynamics (Friedmann $+$ Klein--Gordon with
$V(\varphi) = A\varphi^4(2\ln(\varphi/v_S) - \tfrac{25}{6})$, $A \approx 3.6\times 10^{-6}$,
$\varphi_0 \approx 22.9\,v_S$, $\varphi_{\mathrm{crit}} \approx 8.0\,v_S$,
$m_\varphi \approx 0.27\,H_0$) gives recollapse at $a_{\max} \approx 8\times 10^{17}$
($\sim 41$ $e$-folds, $t \sim 170/H_0$), whereas $\OL = 0.685$ would require
$a_{\max} \approx 6.7$ --- seventeen orders of magnitude too late. The light scalaron rolls too
slowly from $\varphi_0 \gg \varphi_{\mathrm{crit}}$, so the recollapse lies far beyond
observational horizons (the explicit caveat of Section~\ref{sec:theorem64}) and cannot truncate
the shell lifetime in the relevant range.

\paragraph{The measure lever.} A second attempt replaces the uniform-in-$y$ prior by the
Hartle--Hawking amplitude $|\Psi(\sigma_0)|^2 \propto e^{-8\ag\sigma_0}$, the canonical
no-boundary weight over cosmological histories. The amplitude alone does not fix the observer
weighting; combined with different companion measures it ranges over the entire admissible
interval --- $\text{HH}\times dN \to 0.00$, $\text{HH}\times dt \to 0.04$,
$\text{HH}\times V_4 \to 0.78$, $\text{HH}\times(\text{phase space}) \to 1.00$. The framework's
own canonical choice (the amplitude \emph{is} the $V_4$ weighting) gives $0.78$, not $0.685$. The
companion measure one would need is physically the observer density --- when and where observers
exist --- which no axiom fixes. This is the \emph{anthropic} measure problem of quantum
cosmology, and in that form it is categorical, not technical: weighting a history ensemble by an
observer density that the theory does not supply, $\OL$ cannot be pinned by any
prior-independent rule.

\paragraph{The causal route avoids the measure problem entirely.} The anthropic difficulty arises
only because one weights an \emph{ensemble} of histories. A causal variational closure removes the
ensemble altogether: fixing $\Lambda$ by stationarity over a single past light cone
(Section~\ref{sec:causal}) needs no observer measure. The self-consistency
$\Lambda = 3\OL H_0^2 \sim c/t_U^2$ then has a prior-free solution depending only on the
dimensionless light-cone coefficient $c = A_{\partial M}/(t_U^{-1}V_M) = O(1)$, where $A_{\partial
M}$ is the boundary $3$-volume and $V_M$ the four-volume of the past light cone. Two questions
must be separated: which boundary hypersurface, and which functional form.

\emph{The hypersurface is forced.} It is not a free choice: the maximal-enclosed-$3$-volume
slice (the ``belly'') coincides \emph{exactly} with the apparent horizon, the covariant
$\theta=0$ surface, since $\dot r = Hr-1 = 0 \Leftrightarrow r = 1/H$ for the proper light-cone
radius $r = a\chi_{\rm lc}$. Two independent characterizations --- maximal $3$-volume and the
covariant Bousso/Hayward $\theta=0$ surface --- select the same slice (numerically $r H = 1.000$).
The particle horizon, a candidate alternative, is the big-bang tip with $rH\approx 55$ and is
\emph{not} a covariant causal boundary; it is excluded. The apparent-horizon choice is therefore
forced, and the earlier impression of a $\sim\!\times 500$ ``definition dependence'' was an
artifact of comparing the genuine causal boundary with a non-boundary.

\emph{The functional form is not forced.} With the unique apparent-horizon surface, no
available principle fixes the dimensionless closure. The natural candidate to \emph{derive} it
--- KP stationarity, $\Lambda = \tfrac14\avg{R}$ over the averaging $4$-volume, here restricted
to the causal diamond, $\Lambda = \tfrac14\avg{R}_M$ over $V_M$ --- \emph{fails}: the past-cone
average is dominated by the early-matter $a^{-3}$ curvature, giving $\avg{R}_M/4 \sim 14 \gg 3\OL$
with no self-consistent fixed point. Apparent-horizon thermodynamics do not help either: the
first law at the horizon (Cai--Kim) reproduces the Friedmann equation identically and imposes no
$\Lambda$ constraint, while de~Sitter entropy saturation gives the trivial $\OL=1$. And the bare
dimensional closures spread across almost the whole admissible range: over the unique
apparent-horizon surface, natural choices give $\OL \in \{0.489,\,0.725,\,0.929,\,0.938,\,0.959\}$
(median $0.93$), with no cluster near the observed value. The often-quoted
$3\OL = V_3({\rm AH})/(V_M t_U) \to 0.725$ is therefore \emph{one selective choice}, not a
distinguished one --- most natural closures land near $0.93$. (The earlier band value $c=1.86$ was
\emph{circular}, choosing $L=t_U$ and reading $c=\Lambda t_U^2$ off the input, and is superseded;
Shaw--Barrow's sharp $\Ok=-0.0056$ likewise needs extra physical inputs, the QCD condensate
$\zeta_b\approx\tfrac12$ and the inflationary $e$-fold count $N$, not pure geometry.) So the
causal route does \emph{not} predict $0.685$, nor even a narrow band: it forces the hypersurface
(apparent horizon) and an $O(1)$ range, but the closure functional is undetermined and the
obvious candidates (KP stationarity, horizon thermodynamics) are ruled out. The open question is
thereby sharpened --- from ``a free prior'' through ``which boundary'' (answered: the apparent
horizon) to ``which closure functional,'' which remains genuinely open. What is fixed prior-free
is therefore: the dimensionless ratio $\Ok/\OL = \ag/(8\pi^2)$ (theorem), the magnitude and
coincidence to order unity, the forced apparent-horizon hypersurface, and an $O(1)$ range for
$\OL$ --- not the sharp value, which is pending a derived closure functional, not a free measure.

\section{The sign: anomaly-forced curvature and the Janus point}
\label{sec:sign}

The magnitude of $\Ok$ follows from Eq.~\eqref{eq:central}; the sign is fixed by the cap action
itself. Integrating the stationarity condition $d\Gamma/d\theta_0 = 3\ag\sin^3\theta_0 +
3\cot\theta_0$ gives the reduced cap action $\Gamma(\dtheta) = 3\ag\sin\dtheta - \ag\sin^3\dtheta
+ 3\ln\cos\dtheta$, whose only real stationary point in $(-\tfrac\pi2,\tfrac\pi2)$ is
$\dtheta^* = +0.16391$ for the physical $\ag = +31/180 > 0$ --- an \emph{open} universe
($\Ok>0$). There is no closed-universe saddle: a negative $\dtheta^*$ would require $\ag<0$. The
saddle \emph{location} depends only on the action, whose anomaly source $+\ag$ is the
prescription-independent photon heat-kernel coefficient; the Hartle--Hawking versus Vilenkin
choice ($\Psi \propto e^{-S_E}$ versus $e^{+S_E}$) reweights this same saddle but does not move
it. The sign is therefore forced by $\mathrm{sign}(\ag) > 0$ together with the dominance of the
real cap saddle (B1, $n_{\mathrm{complex}}=0$), not by a branch postulate. The Wick
rotation is the established Bucher--Goldhaber--Turok open-inflation construction
\cite{BGT1995,HawkingTurok1998}: the contour $\theta = \pi/2 + i\tau$ carries the compact cap
into a hyperbolic foliation with the conical singularity off the integration contour, so the
smooth-beginning obstruction of \cite{Feldbrugge2017} does not apply. The Janus point --- the
moment of vanishing extrinsic curvature \cite{Barbour2014} --- is the equatorial $S^3$ where the
two hemispheres meet under a $Z_2$ reflection.

\paragraph{Status of the sign.} The curvature sign separates into two distinct sectors. The
\emph{curvature sector} --- whether the cap deforms toward an open or a closed continuation ---
is fixed: the cap action has a unique real saddle at $\dtheta^* > 0$ for $\ag>0$, and the
explicit computation (Section~\ref{sec:found}) confirms its location is independent of the
$e^{\mp S_E}$ weighting, so the open sign is forced by $\mathrm{sign}(\ag)>0$ and the real-saddle
dominance B1, modulo the prescription-independence of the anomaly-induced action (standard for a
matter heat-kernel coefficient). The \emph{fluctuation-stability sector} --- whether the
Lorentzian continuation of the lapse integral is well defined and normalizable --- is the
Feldbrugge--Lehners--Turok question: the explicit lapse-saddle count is boundary-condition
dependent (four saddles under Dirichlet data, two under Neumann), so the well-definedness of the
Lorentzian no-boundary measure is a foundational input of quantum cosmology, shared by every
no-boundary model and not specific to this framework. The two sectors are independent: the FLT
ambiguity concerns the stability of fluctuations, not the sign of the spatial curvature, which is
already fixed in the Euclidean cap by $\ag>0$. Empirically, DESI DR2 places the open prediction
at $0.74\sigma$ and disfavours the closed (Vilenkin) reading at $4.6\sigma$.

\section{Foundational computations: sign, $B1'$, and the reference epoch}
\label{sec:found}
This section records explicit computations targeting the three items the framework had carried
as non-theorem-level. The outcome is
mixed and stated without embellishment: one item is closed, one is reclassified and closed in its
physical part, and one is confirmed irreducible.

\paragraph{Curvature sign (closed).} The reduced cap action $\Gamma(\dtheta) = 3\ag\sin\dtheta -
\ag\sin^3\dtheta + 3\ln\cos\dtheta$, obtained by integrating the stationarity condition of
Section~\ref{sec:sign}, has exactly one real stationary point in $(-\tfrac\pi2,\tfrac\pi2)$:
$\dtheta^* = +0.16391$ for $\ag = +31/180$, with $\Gamma(\dtheta^*) = 0.0431$. For $\ag<0$ the
unique saddle moves to $\dtheta^* = -0.16391$. Because the saddle location is determined by
$\Gamma$ alone --- whose anomaly source $+\ag$ is the prescription-independent photon
heat-kernel coefficient --- the Hartle--Hawking versus Vilenkin choice ($e^{-S_E}$ versus
$e^{+S_E}$) reweights but does not relocate it. Hence the open sign is forced by
$\mathrm{sign}(\ag)>0$ and the real-saddle dominance B1, not by a branch postulate.

\paragraph{$B1'$: the scalar embedding (cap dominance preserved).} Two layers establish that the
cap-saddle dominance survives the enlargement beyond the one-dimensional reduced action. First,
at one loop the deformed cap $D^4(\theta_0=\pi/2+\dtheta^*)$ is a polar cap on $S^4$ with $S^3$
slices, so $SO(4)$ is unbroken, $\ell$ remains a good quantum number, and the fluctuation Hessian
is exactly block-diagonal in $(\ell,\text{sector})$ with no $\ell$--$\ell'$ mixing; its positivity
is established analytically by domain monotonicity (Section~\ref{sec:cap}), so the cap is a
Morse-index-$1$ point and $n_{\mathrm{complex}}=0$ is preserved exactly.

Second --- the perturbative embedding proper --- the scalar matter field $\varphi$ is restored as
an independent degree of freedom, giving the two-dimensional action $\Gamma_{2\mathrm{D}}(z,\varphi)
= \Gamma_{\mathrm{geom}}(z) - \tfrac12\mu_{\mathrm{eff}}^2\varphi^2 V_{\mathrm{cap}}(z)$ with
$\Gamma_{\mathrm{geom}}(z) = 3\ag(\sin z - \tfrac13\sin^3 z) + 3\ln\cos z$ (identical to the cap
action of Section~\ref{sec:sign}) and $V_{\mathrm{cap}}(z)=\int_0^z\sin^3 t\,dt$. The $1$D saddles
lift trivially: since the Coleman--Weinberg potential satisfies $V_{\mathrm{CW}}(0)=0$, one has
$\Gamma_{2\mathrm{D}}(z,0)=\Gamma_{\mathrm{geom}}(z)=\Gamma_{1\mathrm{D}}(z)$ exactly, so the eight
$1$D saddles are also $2$D saddles with $\mathrm{Im}\,\Gamma$ unchanged (Lemma B1-Stokes inherits).
The scalaron-vacuum saddle at $\varphi=v_S e^{11/6}$ lies only $\sim A_{\mathrm{M5}}\approx
1.2\times10^{-6}$ below the geometric cap action (a relative shift $\sim3\times10^{-5}$ of
$\Gamma_{\mathrm{geom}}\approx0.043$), and the single additional $2$D-specific saddle, at
$z\approx-0.656+1.241i$, has $\mathrm{Im}\,\Gamma\approx1.87\neq0$ --- no Stokes line connects it
to the real cap. Cap dominance is therefore preserved under the scalar perturbation, with the
correction suppressed by $A_{\mathrm{M5}}$.

\emph{Complete thimble enumeration.} The enumeration can be closed rigorously rather than by
sampling. Substituting $u=e^{iz}$ turns the saddle equation $\ag\cos^4 z = \sin z$ into the
polynomial $\ag(u^2+1)^4 + 8i(u^5-u^3)=0$, of degree $8$; by the fundamental theorem of algebra
this returns \emph{all} eight saddles exactly. One is the real cap ($z^*=0.1639$,
$\mathrm{Im}\,\Gamma=0$); the other seven are complex with $|\mathrm{Im}\,\Gamma|\ge1.87>0$. On
the scalaron branch $\varphi=v_S e^{11/6}$ the $\varphi$-equation decouples and the $z$-equation
becomes the $\kappa$-perturbed degree-$8$ polynomial ($\kappa=\mu_{\mathrm{eff}}^2\varphi^2/6
\approx2.4\times10^{-3}$), again exactly eight saddles --- one near-real (the scalaron, at
$\mathrm{Re}\,\Gamma$ lower than the cap by $\sim A_{\mathrm{M5}}$) and seven complex with
$|\mathrm{Im}\,\Gamma|>0$. The decisive step is Picard--Lefschetz: along any Lefschetz thimble
$\mathrm{Im}\,\Gamma$ is constant. The real no-boundary contour carries $\mathrm{Im}\,\Gamma=0$,
while every one of the fourteen complex saddles has $\mathrm{Im}\,\Gamma\neq0$, so each complex
thimble keeps $\mathrm{Im}\,\Gamma=\text{const}\neq0$ and never reaches the real axis. Hence
$n_{\mathrm{complex}}=0$ \emph{rigorously} --- the conclusion follows from the constancy of
$\mathrm{Im}\,\Gamma$ along thimbles, not from a finite numerical saddle scan. The dominant contribution is the pair of near-degenerate real saddles (cap and
scalaron, separated by $\sim A_{\mathrm{M5}}\sim10^{-6}$), and the cap sector dominates the
no-boundary wave function. The only ingredient outside the polynomial closure is the set of
complex log-branches of $V'(\varphi)=0$ ($\varphi^2=v_S^2 e^{11/3+2\pi i n}$): these are
renormalization artifacts of the Coleman--Weinberg logarithm rather than saddles of the
fundamental theory, and are in any case exponentially suppressed, so they do not affect
$n_{\mathrm{complex}}=0$.

\paragraph{Embedding into the full no-boundary path integral (semiclassical).} The pieces above
assemble into a complete semiclassical evaluation of
$\Psi_{\mathrm{NB}} = \int\!\mathcal{D}g\,\mathcal{D}\varphi\;e^{-S_E[g,\varphi]}
= \sum_a n_a\, e^{-\Gamma_a}\,[\det{}' M_a]^{-1/2}\,(1+O(\hbar))$.
(i)~The dominant saddle is the maximally symmetric cap $D^4\subset S^4$, so the saddle search
reduces to the homogeneous mini-superspace $(z,\varphi)$ whose sixteen saddles are completely
enumerated above; Picard--Lefschetz gives $n_{\mathrm{complex}}=0$, leaving the two real saddles
(cap and scalaron, gap $\sim A_{\mathrm{M5}}$) with the cap dominant. (ii)~The inhomogeneous
($\ell\ge1$) modes are Gaussian about the cap with a positive-definite Hessian (domain
monotonicity, Section~\ref{sec:cap}); their $\zeta$-regularized determinant carries the scale
anomaly $\zeta(0;\text{Maxwell};S^4)=-\tfrac{31}{45}=-4\ag$, so the one-loop determinant
$[\det{}'M]^{-1/2}$ is itself the carrier of the prediction $\ag/(8\pi^2)$. The inhomogeneous
modes have vanishing overlap with the anomaly source on $D^4$ by parity (the $\ell=1$ $S^3$
harmonic integrates to zero against $E_4$), so the factorization homogeneous~$\times$~inhomogeneous
is exact at quadratic order. (iii)~By the Adler--Bardeen theorem the Type-A coefficient $\ag$
receives no higher-loop corrections, so the leading prediction is protected; the residual
higher-loop shift is of order $A_{\mathrm{M5}}\sim5\times10^{-4}$ relatively, below DESI~DR2
sensitivity. (iv)~The single $\ell=0$ conformal/lapse mode is rotated to convergence by
Gibbons--Hawking--Perry. The honest limit is precise: this is \emph{one-loop exact} for the
Type-A prediction, not an all-orders non-perturbative evaluation. Full convergence of the
Euclidean quantum-gravity path integral remains a generic open problem (the conformal-factor
issue, mitigated but not resolved by GHP), shared with every approach to semiclassical quantum
gravity and not claimed here.

\paragraph{All-orders stability of the prediction.} The convergence question splits cleanly, and
the physically relevant half is provable. The observable is the ratio $R=\ag/(8\pi^2)$, and $R$
is the quotient of two $\hbar$-independent quantities, so it receives no perturbative correction
at any loop order. First, CSG treats gravity classically: graviton loops do not exist, so the
genuine obstruction to convergence of a Euclidean gravity path integral --- the unboundedness of
the conformal-factor action and the perturbative non-renormalizability of Einstein gravity, both
graviton-loop effects --- is structurally absent, and the remaining integral is matter on a
classical saddle. Second, on the round $S^4$ the Weyl tensor vanishes ($W^2=0$), so Type-B drops
out; the surviving Type-A coefficient $\ag=31/180$ is Adler--Bardeen one-loop exact, receiving no
higher-loop correction. Third, $8\pi^2=\int_{D^4}Q_4$ is a topological charge, a
discrete/rational invariant independent of $\hbar$ to all orders. Fourth, the saddle sum is
finite-dimensional with $n_{\mathrm{complex}}=0$. Hence $R$ is all-orders well-defined within CSG.
The honest boundary is equally precise: the absolute Borel summability of the four-dimensional
$\varphi^4$ matter series (the generic triviality/Landau question) and the normalization of
$|\Psi|^2$ are \emph{not} settled here, but both cancel from the ratio and are protected by its
topological character --- they do not bear on $R$. This is not a solution of the general
quantum-gravity path integral; it is the decoupling of the prediction from that open problem, and
it holds within the CSG assumption of classical/conformal gravity (itself an empirical question).

\paragraph{Fluctuation stability in the Euclidean sector (resolved).} The single negative
direction of the cap saddle is the conformal/scale-factor mode: $d^2\Gamma/d\dtheta^2(\dtheta^*) =
-3.33 < 0$, while all transverse modes are positive (B1$'$ above). This is the Euclidean
Gibbons--Hawking--Perry conformal-factor situation, and its standard resolution applies: rotating
the conformal mode along the imaginary axis, $\dtheta = \dtheta^* + is$, makes $\mathrm{Re}\,
\Gamma$ grow ($0.043 \to 0.060 \to 0.109$ for $s = 0, 0.1, 0.2$), so $e^{-\Gamma}$ decays and the
Gaussian over the rotated mode converges. The Euclidean one-loop no-boundary wavefunction is
therefore well defined. The Feldbrugge--Lehners--Turok obstruction concerns the \emph{Lorentzian}
lapse integral --- a different definition of the path integral, which this framework does not
adopt; the Euclidean cap is its primary object. Whether the Euclidean (GHP-rotated) or the
Lorentzian definition is the physically correct quantum gravity is an open foundational debate,
not a CSG--KP-specific gap.

\paragraph{The reference epoch A5c (no fit freedom, identification irreducible).} The physical
density ratio $\Ok(z)/\OL(z) \propto (1+z)^2$ is epoch-dependent (today $2.18\times10^{-3}$, at
$z=1100$ about $2.6\times10^3$), so a reference epoch must be specified. It introduces no fit
freedom: given the measured $\OL$, the prediction $\Ok = (\ag/8\pi^2)\OL$ is unique, with no
tunable parameter. A5c is the identification rule ``the Euclidean topological number equals the
observable at the observation epoch,'' forced by testability (we measure at $z=0$). Anchoring
instead at asymptotic de~Sitter ($\OL\to1$) would shift the predicted $\Ok$ to $\ag/8\pi^2$, a
$\le46\%$ change; that the matching surface is the present epoch rather than the asymptotic future
is a physical identification, not derivable from first principles.

\paragraph{Status after three iterations.} Under the framework's own (Euclidean,
observation-anchored) definitions, every item closes: the curvature sign by $\ag>0$, $B1'$ at one
loop ($SO(4)$-exact), Euclidean fluctuation stability by GHP rotation plus transverse positivity,
and A5c without fit freedom. The two genuinely irreducible residues are interpretational rather
than computational, and shared with the wider field rather than specific to this framework: the
choice between a Euclidean and a Lorentzian path-integral definition, and the physical
identification A5 of the Euclidean invariant with the observed curvature. These are not gaps to be
computed away; they are the irreducible interpretational content, and observation rather than
further algebra is what bears on them.

\section{Physical interpretation: the Photon-Field Substrate}
\label{sec:psf}

The mathematical structure admits a physical reading, the Photon-Field Substrate (PSF), built
on five postulates: universal vacuum entanglement; the photon as universal massless carrier;
the conformal anomaly as the substrate's residual stress; the Komargodski--Schwimmer $a$-theorem
as the infrared selection principle; and holographic accumulation $N_{\mathrm{eff}} =
\Mpl^2/H^2$. The number $\ag/(8\pi^2)$ then emerges from two independent identities --- the de
Sitter structural identity $\rho_{\mathrm{anom}}/(3H^4) = \ag/(8\pi^2)$ and the hemisphere
topological quotient $\Delta\Gamma/\CQ = \ag/(8\pi^2)$ --- joined by two physical
identifications (the holographic Friedmann density and the curvature observable A5). PSF is
complementary to the CSG--KP mathematical structure in the Bohr sense \cite{Bohr1928}, and is
consistent with locality and with loophole-free Bell tests
\cite{Hensen2015,Giustina2015,Shalm2015}, the correlations being kinematic vacuum entanglement
(Reeh--Schlieder \cite{ReehSchlieder1961}) rather than signalling.

\paragraph{The magnitude hierarchy is carried by P5, and P5 is a postulate.} The topological
ratio $\ag/(8\pi^2) \approx 2.2\times 10^{-3}$ is dimensionless and clean, but its physical
realization carries the full cosmological-constant hierarchy. A single-mode anomaly inserted
naively into the Friedmann balance gives
\begin{equation}
\Omega_{\mathrm{anom}}^{\mathrm{naive}} = \frac{\ag}{8\pi^2}\cdot\frac{H^2}{\Mpl^2}
\sim 10^{-3}\cdot 10^{-122} \sim 10^{-125},
\end{equation}
$122$ orders of magnitude below the observed curvature density. The holographic accumulation
postulate P5, $N_{\mathrm{eff}} = \Mpl^2/H^2$, compensates this exactly:
$\rho_{\mathrm{eff}} = N_{\mathrm{eff}}\,\rho_{\mathrm{anom}} = (3\ag/8\pi^2)\,\Mpl^2 H^2$, so
$\Omega_{\mathrm{eff}} = \ag/(8\pi^2)$. The same hierarchy appears on three levels --- $O(1)$ (the
dimensionless ratio), $10^{122}$ (energy density, $\Mpl^2/H^2$, equal to the horizon area $A_2$ in
Planck units), and $10^{244}$ (four-volume, $(\Mpl/H)^4 = V_4$) --- and the $10^{240}$ failure mode
of the dynamical Mechanism~2 (Section~\ref{sec:uniqueness}) is exactly this $V_4$ factor, not an
artifact. P5 is what bridges the dimensionless ratio to the observable density. It is supported by
the 't~Hooft--Susskind area bound (which $N_{\mathrm{eff}}$ saturates up to a factor $\pi$) and by
the dimensional identity $V_4 = A_2^2$ (algebraically trivial, a consistency check only). The
substantive content is not the count but its \emph{coherence}: the modes must add linearly
($\rho_{\mathrm{eff}} = N\rho$), not in random phase ($\sqrt N\rho$, which would give
$\Omega \sim 10^{-64}$, catastrophically wrong). Three structural mechanisms enforce coherence:
(K1) Reeh--Schlieder vacuum entanglement makes the local trace expectation
$\avg{T^\mu{}_\mu} = \ag E_4$ pointwise deterministic rather than stochastic; (K2) the integrated
anomaly is topologically fixed, $\int_{S^4} E_4 = 64\pi^2$, independent of the mode number, so the
modes must sum linearly to reproduce it; (K3) the KP global constraint
$\Lambda V_4 = \int\sqrt g\,\avg{T^\mu{}_\mu}$ synchronizes sub-regions, forbidding out-of-phase
oscillation. Each independently excludes the random-phase regime. The \emph{coherence} half can
moreover be reduced beyond bare motivation: in Mottola's effective theory of the conformal anomaly
\cite{Mottola2010,Mottola2022}, the anomaly vacuum energy is a \emph{condensate} of a $4$-form field
strength $F=dA$ --- the local conformalon form of the nonlocal anomaly action --- i.e.\ a single
coherent classical amplitude over the horizon volume, and a condensate adds in phase ($\times N$) by
construction, the $\sqrt N$ regime requiring incoherent modes a condensate does not possess. The
infrared-relevance of the same action places the backreaction at the horizon scale rather than the
UV, with the value set by horizon boundary conditions \cite{Mottola2010}; Mottola's
$\Lambda\to0$ infrared fixed point is then the sequestering of the bare term
(Section~\ref{sec:sequester}), while the finite $4\ag$ is the anomaly-residual condensate value ---
two distinct contributions, consistent. What is \emph{still} not available is a
strict Schwinger--Keldysh or BV-master-equation derivation fixing the \emph{exact} $O(1)$
normalization of $N_{\mathrm{eff}}=(\Mpl/H)^2$, and Mottola's conformalon EFT, though established, is
not universally accepted. Thus P5 is the single quantitatively decisive
element of the framework that is not derived from first principles: it carries the entire
$122$-order magnitude hierarchy, with its coherence now a \emph{consequence of the anomaly
condensate} rather than a bare postulate, leaving the residual open piece as the one $O(1)$ horizon
normalization. The Euclidean ratio is theorem-level; the magnitude of the physical curvature
density rests on P5.

\paragraph{The normalization Mottola leaves open is fixed by the Q-charge; P5 is the declared
founding assumption.} Mottola's mechanism does not, by itself, fix the value: it ``depends on
macroscopic boundary conditions at the cosmological horizon scale'' \cite{Mottola2010}, and the
conformal theory --- hence $\ag$ --- is unspecified. These are exactly what CSG--KP supplies and a
generic anomaly EFT does not. The boundary condition is the $Q$-curvature charge of the hemisphere,
$\CQ=8\pi^2$ (Chang--Yang), and the conformal theory is the photon ($\ag=31/180$), so the
accumulation is the de Sitter entropy per $Q$-charge,
\begin{equation}
N_{\mathrm{eff}} = \frac{S_{dS}}{\CQ} = \frac{8\pi^2(\Mpl/H)^2}{8\pi^2} = \left(\frac{\Mpl}{H}\right)^2 ,
\end{equation}
with the $O(1)$ factor \emph{exactly} unity because $S_{dS}$ and $\CQ$ share the same $8\pi^2$
normalization. The finite value Mottola's $\Lambda\to0$ fixed point does not provide is then the
anomaly residual $\OL=4\ag$, fixed by $\ag$ and $\CQ$. What was an undetermined $O(1)$ in the
generic EFT is therefore closed by the framework's own topology. The residual is not a hidden gap
but a single \emph{declared} founding assumption: P5 $:=$ ($N_{\mathrm{eff}}=S_{dS}/\CQ$, coherent
accumulation as entropy per $Q$-charge), as fundamental as the holographic principle itself. Given
P5, the framework is internally complete --- value, sign, ratio, and dark-sector interaction all
follow --- and its honesty lies in the explicitness of that one postulate together with a sharp
falsification channel (the curvature ratio, DESI/Euclid). This is the ordinary structure of a
physical theory, not an outstanding deficit.

\paragraph{P5 is built from established results: the entanglement area law and Jacobson's
equilibrium, exact for the photon.} P5 has two parts, and the literature supplies both (the value
excepted), so it is the application of recognized results to the photon substrate rather than an
isolated postulate. First, its area scaling $N_{\mathrm{eff}}=(\Mpl/H)^2$ \emph{is} the entanglement
area law: the vacuum entanglement entropy of a region scales with its boundary area as $A/4L_p^2$,
the Bekenstein--Hawking form \cite{Jacobson2016,RyuTakayanagi2006}, and $S_{dS}/\CQ=(\Mpl/H)^2$.
Second, the vacuum-entanglement-to-$\Lambda$ link is Jacobson's entanglement equilibrium
\cite{Jacobson2016}: the semiclassical Einstein equation with a cosmological constant holds if and
only if the vacuum entanglement entropy of small causal diamonds is stationary at fixed volume ---
and crucially this result is \emph{exact} for first-order variations of the vacuum of \emph{conformal}
fields (only modulo a conjecture for nonconformal ones). The photon is conformal and the unique
infrared-surviving field (Komargodski--Schwimmer), so Jacobson's theorem applies to the CSG photon
substrate exactly: this lifts K1 (Reeh--Schlieder vacuum entanglement) from a structural motivation
to an established theorem for precisely the relevant field type. Moreover the maximally symmetric
causal diamond that carries the equilibrium is itself de Sitter space with a cosmological constant
\cite{JacobsonVisser2019}, i.e.\ the CSG cap, and the negative diamond temperature it requires
matches the Hartle--Hawking weighting $e^{-S_E}$; while the coherence ($\times N$, not $\sqrt N$) is
exactly the strict area law that deriving the Einstein equation demands \cite{Verlinde2017}. What no
such programme provides is the \emph{value}: Jacobson leaves $\Lambda$ an integration constant and
Verlinde's dark-energy scale carries contested assumptions. The value $\OL=4\ag$ is fixed by CSG's
$\ag$ (photon) and $\CQ$ (hemisphere); an assumption-free theorem \emph{with} the value is absent
across the literature, a gap of the field rather than of CSG specifically. P5 thus stands on the
area law and Jacobson's equilibrium for the conformal photon, with the value as CSG's own input.

\paragraph{P5 as Jacobson's theorem with two premises derived and the third triangulated.}
Jacobson's derivation rests on three inputs, and the decisive observation is that CSG's anomaly
and infrared structure supply precisely the two that are contested in the literature. \emph{(i)} A
universal area density of vacuum entanglement entropy: by Solodukhin and Casini--Huerta the
universal four-dimensional entanglement entropy \emph{is} the conformal anomaly, its type-A part
fixed by the topology of the entangling surface (the extrinsic-curvature contribution drops for the
spherical de Sitter hemisphere), so the universal coefficient is $\ag$, rendered finite by the
topological charge $\CQ$ rather than by a cutoff. The known Maxwell subtlety --- the naive free
logarithmic coefficient differs from the anomaly because of edge modes --- is resolved through the
four-sphere partition function, which is exactly the object $\zeta(0;\text{Maxwell};S^4)=-4\ag$ that
CSG uses. \emph{(ii)} A conjecture that the nonconformal variation matches the conformal form, which
Casini, Galante and Myers found to conflict for relevant operators of low conformal dimension: this
case is voided in CSG because the Komargodski--Schwimmer $a$-theorem leaves only the conformal photon
in the infrared, so the problematic nonconformal sector never enters. The remaining input,
\emph{(iii)} that the vacuum is an entanglement equilibrium (the entropy stationary at fixed volume),
is the one CSG does not prove; it is concretized as the cap-saddle stationarity
$\mathrm{d}\Gamma/\mathrm{d}\theta=0$ at the de Sitter diamond. This residual premise is not proved
but \emph{triangulated} by convergent results of the full pipeline: the cap saddle $\delta\theta^*$ is
a genuine action stationary point; the stationary action equals the anomaly invariant
$\int_{S^4}\langle T\rangle=4\ag$; the Riegert scale flow $\mathrm{d}\Gamma/\mathrm{d}\sigma=4\ag$
reaches the same invariant by an independent route; the entropy whose stationarity is at issue is
$N_{\mathrm{eff}}=S_{dS}/\CQ$; and the KP volume constraint $\Lambda V_4=\int\!\sqrt{g}\,\langle
T\rangle$ is the Jacobson--Visser first law of the diamond. All five identify the cap saddle with the
stationarity of the anomaly action and entropy at the de Sitter diamond, which Jacobson--Visser
reformulate as the entanglement equilibrium \cite{JacobsonVisser2019}; the numerical convergence of
the action invariant and the Riegert flow on the same $4\ag$ by independent computations is the
sharpest leg of the triangulation. P5 is thereby lifted from an isolated postulate to Jacobson's
entanglement-equilibrium theorem for the conformal photon, with premises (i) and (ii) derived from
CSG and (iii) triangulated. This is a qualitative reduction, not an assumption-free closure: premise
(iii) remains a physical equilibrium hypothesis --- Jacobson's own, and broadly accepted --- the
Maxwell edge-mode coefficient is still under debate, and Jacobson's equivalence is established at
first order in the variations.

\paragraph{The honest boundary, drawn smaller: one equilibrium premise.} Once the pipeline is
convergent, two reductions follow. First, the Hartle--Hawking versus Vilenkin branch choice ceases
to be an independent postulate. The entanglement-equilibrium premise is defined about the maximally
symmetric vacuum, which for the conformal field is the Bunch--Davies state; by Gibbons--Hawking this
state restricted to the static patch is a thermal equilibrium, and the Hartle--Hawking vacuum is
stationary, $\rho\sim e^{-\beta H}$, whereas Vilenkin's tunnelling state is an out-of-equilibrium
initial condition. Given the equilibrium premise the branch is therefore forced to Hartle--Hawking,
with the caveat that in the bare wave-function formalism the sign is not fixed absolutely, so the
forcing holds only relative to the premise. Second, the seemingly separate assumptions --- A5
(the anomaly action $4\ag$ is the curvature), premise (iii) (the vacuum is an entanglement
equilibrium), and P5 ($N_{\mathrm{eff}}=S_{dS}/\CQ$) --- are one statement in three forms, since
$\OL=4\ag$ is exactly equivalent to the normalization $N_{\mathrm{eff}}=S_{dS}/\CQ$ and Jacobson
identifies the equilibrium stationarity with that same action invariant. Counting them separately
triple-counts a single premise. The exhaustion no-go reinforces this: no dynamical route yields
$4\ag$ --- every dynamics gives a different number --- and only the action invariant
$\int_{S^4}\langle T\rangle=4\ag$ does, so the value cannot be dynamical and the governing principle
must be one of stationarity, which is precisely the equilibrium premise. The ledger thus reduces from
two postulates to one: a single equilibrium premise --- Jacobson's own, broadly accepted, and forced
in form by the no-go --- from which the branch follows and from which, with premises (i) and (ii)
supplied by CSG, the value follows. The residual is the strict functional identity between the
cap-saddle stationarity and Jacobson's entanglement stationarity, triangulated by the convergent
pipeline but not proved. This is the smallest honest boundary the framework reaches: one equilibrium
postulate and one triangulated identity, not a first-principles closure.

\paragraph{Reverse anchoring: CSG--KP supplies the number established frameworks must postulate.}
P5 and the absolute identification do not stand alone; they connect to established holographic
dark-energy programmes by a reverse bootstrap. Those programmes fix the dark-energy density
through a holographic relation containing one dimensionless $O(1)$ number they must \emph{fit or
postulate}; CSG--KP supplies that number from $\ag$, and in return the programme supplies the
energy-\emph{density} form that the no-go of Section~\ref{sec:ladder} showed the anomaly cannot
produce dynamically. Concretely, holographic dark energy \cite{CohenKaplanNelson1999,Li2004}
posits the genuine density $\rho_{\mathrm{DE}} = 3c^2\Mpl^2/L^2$ with $c\sim0.8$ fitted; taking the
Hubble-scale cutoff $L=1/H$ and the CSG--KP value $c^2 = 4\ag$ gives
\begin{equation}
\rho_{\mathrm{DE}} = 3(4\ag)\Mpl^2 H^2 \quad\Longrightarrow\quad \OL = 4\ag ,
\end{equation}
so $4\ag$ appears as a real density fraction, not merely as an action invariant, and the
predicted coefficient $c=\sqrt{4\ag}=0.830$ matches the holographic fit. The same number recurs as
the holographic-equipartition factor: the bridge ratio $\rho_\Lambda/\rho_{\mathrm{anom}} =
32\pi^2(\Mpl/H)^2 = 4\,S_{dS}$ is the de~Sitter entropy times $\chi(S^4)^2=4$, with the surface
form $N_{\mathrm{sur}}=N_{\mathrm{bulk}}$ supplied by emergent gravity \cite{PadmanabhanCosMIn2013}.
\emph{Honest caveats.} The Hubble-cutoff HDE has $w_{\mathrm{DE}}=0$ classically \cite{Hsu2004};
recovering $w_{\mathrm{eff}}=-1$ together with a constant $\rho_{\mathrm{DE}}/\rho_m$ (the
coincidence) requires a dark-sector interaction \cite{PavonZimdahl2005}. That interaction is not a
free addition, however: the next paragraph shows it is \emph{forced} by the Bianchi identity once
the vacuum runs, with its strength set by $\ag$. And CosMIn \cite{PadmanabhanCosMIn2013} does not
reverse-extract: it fixes the \emph{absolute} scale $\Lambda L_P^2$ through $H_0$, whereas $4\ag$ is
the \emph{fraction}; the two are complementary, not reducible. The net gain is real but bounded:
$4\ag$ now carries a density realization and is anchored independently in the HDE coefficient, the
equipartition factor, and the anomaly action --- no single one proving it, but together removing
its isolation. The residual open step is no longer ``action or density'' but the choice of the
Hubble cutoff and the exact strength of the (Bianchi-forced) interaction.

\paragraph{The dark-sector interaction is forced, not free: running vacuum from the anomaly.}
The interaction the previous paragraph needs is not an extra ingredient but a consequence of energy
conservation. If the vacuum energy runs, $\rho_{\mathrm{vac}}=\rho_{\mathrm{vac}}(H)$, the total
Bianchi identity $\nabla_\mu T^{\mu\nu}_{\mathrm{tot}}=0$ \emph{forces} an exchange with matter
that cannot be switched off; this is the running-vacuum mechanism
\cite{BilicEtAl2007,SolaPeracaula2021}. The forced continuity equation is
\begin{equation}
\dot\rho_m + 3H\rho_m = +3\nu H\rho_m \quad\Longrightarrow\quad \rho_m \propto a^{-3(1-\nu)} ,
\end{equation}
with a single dimensionless strength $\nu$ equal to the $\beta$-function coefficient of the running
cosmological constant --- a QFT quantity, $\nu_{\mathrm{eff}}=\nu_s+\nu_f+\nu_v$, to which vector
bosons contribute even in the conformal limit ($\nu_s=0$ but $\nu_v\neq0$)
\cite{SolaPeracaula2021}. In CSG--KP the only infrared-surviving anomaly source is the photon, so
$\nu=\nu_v(\text{photon})$, and the natural anomaly magnitude is
\begin{equation}
\nu = \frac{\ag}{8\pi^2} = \frac{31}{1440\pi^2} \approx 2.18\times10^{-3} ,
\end{equation}
which lies inside the running-vacuum data-fit band $|\nu|\sim10^{-3}$--$10^{-2}$ \emph{and} is
numerically identical to the curvature ratio $\Ok/\OL$ of Eq.~\eqref{eq:central}. The same anomaly
number thus drives both the curvature prediction and the dark-sector exchange rate, so the
coincidence resolution introduces \emph{no} new parameter, and $w_{\mathrm{eff}}=-1$ is recovered
from the running vacuum rather than from a free Hubble-cutoff fluid, bypassing the Hsu obstruction.
\emph{Two honest caveats remain.} First, $\nu=\ag/(8\pi^2)$ is the natural choice with the correct
magnitude, but the exact $\nu_v$ for a massless vector field via adiabatic regularization has not
been computed here --- the magnitude and the consistency stand, the sharp derivation does not.
Second, generic phenomenological Hubble-cutoff interactions are excluded by CMB/LSS data; running
vacuum is a distinct, field-theoretically grounded class, only partly tested, so the specific
CSG--KP value of $\nu$ must still be confronted with structure-formation data. This reduces point~B
from ``needs a new interaction parameter'' to ``needs the anomaly number $\ag/(8\pi^2)$, already
present as the curvature ratio'' --- progress, not closure. \emph{A third caveat must be stated
plainly:} this mechanism yields $w_{\mathrm{eff}}=-1$ as the CSG--KP \emph{prediction}, but DESI's
$w_0w_a$ data prefer a \emph{dynamical} $w$ at ${\sim}4\sigma$ (if real), a deviation of order
$0.25$--$0.6$ that an anomaly number $\nu\sim2\times10^{-3}$ is two orders of magnitude too small to
produce. The running-vacuum term keeps $w_{\mathrm{eff}}=-1$ self-consistently, but it does
\emph{not} explain --- and cannot absorb --- the DESI $w$-preference; that remains the framework's
sharpest open empirical tension, stated as such in the robustness analysis below
(Section~\ref{sec:status}).

\paragraph{The two open steps are one: a massless source forces the accumulation.}
Pressing the interaction calculation harder settles both residual steps at once, by
showing they are the same assumption. The running-vacuum $O(H^2)$ term that acts today,
$\delta\rho_{\mathrm{vac}}\sim\nu_{\mathrm{eff}}\Mpl^2H^2$, originates in field
\emph{masses} (terms $\sim m_i^2/\Mpl^2$; \cite{SolaPeracaula2021}), so a massless photon
gives \emph{no} $O(H^2)$ contribution, $\nu_v(O(H^2))=0$. The photon enters only at
$O(H^4)$ --- the conformal-anomaly density $\rho_{\mathrm{anom}}=3\ag H^4/8\pi^2$ ---
which today, with $H^2/\Mpl^2\sim10^{-122}$, is $122$ orders too small to act as the
$O(H^2)$ interaction. The naive extraction of $\nu$ from a massless field thus fails. It is
\emph{the same} P5 accumulation $N_{\mathrm{eff}}=(\Mpl/H)^2$ that rescues it: lifting the
$O(H^4)$ term gives the Friedmann fraction
\begin{equation}
\frac{N_{\mathrm{eff}}\,\rho_{\mathrm{anom}}}{3\Mpl^2H^2} = \frac{\ag}{8\pi^2} = \frac{\Ok}{\OL},
\end{equation}
exactly the curvature ratio, and the Hubble cutoff $L=1/H$ of the density bridge is that
\emph{same} de Sitter horizon scale, $N_{\mathrm{eff}}=(\Mpl/H)^2=$ the horizon area in
Planck units. Hence the absolute value (which needed the cutoff) and the interaction
strength (which needed the lift) are not two independent assumptions but one: the coherent
holographic accumulation at the de Sitter horizon. The coherence motivation transfers
intact --- the sharpest piece, K2, anchors both to the same topological invariant, since the
value uses $\int_{S^4}E_4=64\pi^2$ and the interaction's vacuum-running coefficient is fixed
by the same flux through the Riegert flow $d\Gamma/d\sigma=(\ag/16\pi^2)\int_{S^4}E_4=4\ag$.
The net effect is a reduction of the open surface from two seemingly separate conjectures to
one --- coherent saturation of the holographic bound, $N_{\mathrm{eff}}=S_{dS}/(8\pi^2)$,
as fundamental as the holographic principle itself and still lacking a strict
Schwinger--Keldysh or BV derivation. A sharper localization of what is open, not a closure.

\paragraph{Observer-timelessness: the ratio is epoch-invariant, the absolute scale need not be.}
The $N_{\mathrm{eff}} = \Mpl^2/H^2$ compensation has a consequence that reframes the coincidence
``why now''. Because $N_{\mathrm{eff}}\sim H^{-2}$ and $\rho_{\mathrm{anom}}\sim H^4$, the factor
$H$ cancels \emph{entirely}: $\Omega_{\mathrm{eff}} = \ag/(8\pi^2)$ holds at every epoch,
independent of $H$. The photon substrate is lightlike and carries no proper time, so the
prediction --- the ratio $\Ok/\OL$ --- is observer- and epoch-invariant, while the absolute $\OL$
carries the scalaron position $\varphi_0$, i.e.\ ``where/when the observer is'', and is in this
reading observer-dependent rather than predicted. The ``why now'' is then dissolved as
observer-selection on the absolute scale, the timeless ratio needing no tuning. This is the
honest counterpart to the geometric ladder of Section~\ref{sec:ladder}: the ladder is the
sharper but \emph{conjectural} alternative, in which even the absolute $\OL = 4\ag$ is topological
(full $S^4$); the observer-timelessness reading is the conservative one, in which only the ratio
is fixed. The two are in tension --- the ladder would make $\OL$ epoch-marked but topologically
determined --- and which holds depends on the open action-identification of
Section~\ref{sec:ladder}.

\section{Higher-order corrections}
\label{sec:corrections}

\subsection{The two-loop coefficient $c_1 = 2\zeta(3)$}
The Coleman--Weinberg expansion of the scalaron potential carries a two-loop coefficient
$c_1 = 2\zeta(3)$, taken as a \emph{literature-supported} input from Mottola--Vaulin
\cite{MottolaVaulin2006}. The primitive period $6\zeta(3)$ of the underlying sunset master
integral is rigorous; the normalization $6\zeta(3) \to 2\zeta(3)$ is the Mottola--Vaulin
photon-mediated value. A naive $\sigma^3$ sunset gives $\tfrac29\zeta(3)$, the
photon-polarization-enhanced version $\tfrac89\zeta(3)$; the remaining factor (a
Wess--Zumino-cocycle contribution) is supplied by the literature value. This is a deliberate and
clearly bounded reliance on a peer-reviewed result; Section~\ref{sec:m1} (M1-$c_1$) confirms the
literature input rather than revising it.

\subsection{The Coleman--Weinberg amplitude}
At the Hubble scale all massive fields are infrared-decoupled, leaving the photon with effective
coupling $g_{\mathrm{eff}}^2 = \ag$. Naive-dimensional-analysis loop counting with two
Wess--Zumino anomaly insertions over two four-dimensional loop integrations gives the M5
amplitude
\begin{equation}
\label{eq:AM5}
A_{\mathrm{M5}} = \frac{\ag^2}{(16\pi^2)^2} \approx 1.19\times 10^{-6},
\end{equation}
and including the two-loop coefficient, $c_1\,A_{\mathrm{M5}} = 2\zeta(3)\,A_{\mathrm{M5}}
\approx 2.86\times 10^{-6}$. The $(16\pi^2)^2$ structure is fixed by the two-loop counting.

\subsection{Convergence of the loop series}
The three-loop sunset master integral
$T_{3\mathrm{SS}} = -3\zeta(5) + 5\zeta(3)^2 - 2\zeta(3)$ gives (symmetry factor $1/144$,
Antoniadis--Mottola vertex $2/3$)
\begin{equation}
c_2 = -\tfrac{1}{108}\zeta(5) - \tfrac{1}{162}\zeta(3) + \tfrac{5}{324}\zeta(3)^2
  \approx 5.3\times 10^{-3} ,
\qquad
|c_2\,(c_1 A_{\mathrm{M5}})^2| \approx 4.3\times 10^{-14} ,
\end{equation}
where the relevant amplitude is the two-loop value $c_1 A_{\mathrm{M5}} \approx 2.86\times 10^{-6}$
of Eq.~\eqref{eq:AM5}, not $A_{\mathrm{M5}}$ alone.
Each loop order is suppressed by $\ag^2/(16\pi^2)^2 \approx 1.2\times 10^{-6}$; the three-loop
contribution is roughly nine orders of magnitude below DESI DR3 sensitivity. This is a
\emph{proven} bound (FORM-verified master integrals, Section~\ref{sec:m1} M1-3L).

\section{Independent computational verification (M1)}
\label{sec:m1}

The framework's key quantitative claims have been independently recomputed with FORM 5.0.0,
SymPy (exact symbolic arithmetic), and SciPy (numerical eigenvalue solvers). No result
contradicts the main text; the open items in particular are \emph{corroborated}, not revised.
The standalone code accompanying this work reproduces each of the following.
\begin{description}[leftmargin=2.4em,nosep,itemsep=2pt]
\item[M1-Q.] The $Q$-charge $\int_{D^4} Q_4 = 8\pi^2$ is verified by exact symbolic arithmetic
($Q_4 = 6$ on the unit $S^4$; $\int_0^{\pi/2}\sin^3 = 2/3$; $6\cdot 2\pi^2\cdot 2/3 = 8\pi^2$),
matching Lemma~\ref{lem:qcharge}.
\item[M1-$\zeta$.] The DtN identity is verified in both bookkeepings of
Section~\ref{sec:two-books}; multi-gauge cross-checks (Lorenz and Coulomb) reproduce the
hemisphere value and return the closed value $-31/45$.
\item[M1-3L.] The three-loop Coleman--Weinberg bound $|c_2 (c_1 A_{\mathrm{M5}})^2| \approx
4.3\times 10^{-14}$ is reproduced from the FORM master integrals.
\item[M1-H.] The cap-saddle stability is reconfirmed: shooting-method eigenvalues in the scalar,
$1$-form, and TT sectors are all positive on the deformed hemisphere (e.g. scalar
$10.0 \to 8.16$, $1$-form $12.71 \to 10.67$, TT $20.0 \to 17.27$ for the lowest modes), with the
single negative lapse direction unchanged.
\item[M1-$\Lambda$.] The trace-zero identity $\Lambda^* = \tfrac14\avg{R} \Leftrightarrow
\avg{T} = 0$ is verified symbolically; the closed-history requirement (finite $V_4$ for $k=+1$,
divergent for $k=0,-1$) and the unique negative Coleman--Weinberg vacuum at
$\varphi_{\min} = v_S e^{11/6}$ are reproduced.
\item[M1-FRG.] The functional-renormalization-group analysis corroborates the assessment that
Interpretation~A (the RG-attractor route for A2) is insufficient --- the dimensionless ratio
overshoots the target by a factor $\sim13$ with no attractor. A2 is instead \emph{numerically
established} in its physically relevant form, Interpretation~B: the iterative fixed point of the
KP self-consistency map (M5 Banach contraction $L\approx0.18$, sign $G'<0$ analytic, $|G'|<2$
bound numerical, Section~\ref{sec:theorem64}). The
retracted RG reading and the iterative fixed point are not in conflict --- they concern flow in
scale versus iteration at fixed cosmological scale --- and the redshift-constancy of the central
ratio is in any case topological, independent of either.
\item[M1-$c_1$.] The literature input $c_1 = 2\zeta(3)$ is confirmed, not revised.
\end{description}

\section{Independent computational verification (M1), in detail}
\label{sec:m1detail}
This section expands Section~\ref{sec:m1}. The recomputation used FORM 5.0.0, SymPy (exact) and
SciPy (numerical); no result contradicts the main text, and the FRG and $c_1$ items corroborate
the existing treatment rather than revising it.
\begin{itemize}[leftmargin=*,nosep]
\item \emph{$Q$-charge}: SymPy reproduces $Q_4(S^4) = \tfrac16(144-108) = 6$,
$\int_0^{\pi/2}\sin^3 = \tfrac23$, $\CQ = 6\cdot2\pi^2\cdot\tfrac23 = 8\pi^2$ exactly; the
cap-shift is a wave-function argument position, so $L_0$ carries no $\dtheta^*$ correction.
\item \emph{$\zeta(0;\mathrm{DtN}) = -1$}: confirmed in multiple gauges; the Esposito
decomposition $\zeta_{\mathrm{PDF}} = -77/180$, $\zeta_{R_1} = -1/20$,
$\zeta_{\mathrm{GM}}+\zeta_{\mathrm{gh}} = -11/30$ sums to $-38/45$, with the BFK route rigorous
and self-contained.
\item \emph{Three-loop bound (FORM)}: $T_{3\mathrm{SS}} = -3\zeta(5) + 5\zeta(3)^2 - 2\zeta(3)$,
$c_2 \approx 5.3\times10^{-3}$, $|c_2 (c_1 A_{\mathrm{M5}})^2| \approx 4.3\times10^{-14}$ (the amplitude is the two-loop value $c_1 A_{\mathrm{M5}}\approx2.86\times10^{-6}$).
\item \emph{Cap-saddle stability} (analytic, by domain monotonicity; confirmed by shooting):
scalar $10.0\to8.16,\,\dots,\,54.0\to48.06$; transverse $1$-form (Weitzenb\"ock-distinct)
$12.71\to10.67,\,\dots,\,56.25\to50.21$; TT $20.0\to17.3,\,\dots,\,59.7\to53.5$; all positive
(each cap value lies strictly above the closed-$S^4$ limit $\ge4,6,12$),
lapse mode Morse index $1$ unchanged.
\item \emph{FRG corroboration}: an independent Reuter--Saueressig flow finds a clean UV fixed
point but the ratio $R = \kappa/\lambda$ overshoots $\ag/(8\pi^2)$ by a factor $13$ at $t=-1$,
with no attractor --- confirming that A2 as an RG attractor fails. A2 is instead numerically established via
Interpretation~B (the iterative fixed point of the KP self-consistency map, M5 Banach
contraction; sign $G'<0$ analytic, $|G'|<2$ numerical); the \emph{topological} $z$-independence of the ratio (the framework's actual
argument) is untouched by either reading.
\item \emph{$c_1 = 2\zeta(3)$}: FORM reproduces the Mottola--Vaulin normalization; the primitive
period $6\zeta(3)$ is rigorous and the factor $6\zeta(3)\to2\zeta(3)$ is the literature input.
\end{itemize}

\section{Observational status and forecasts}
\label{sec:obs}

\subsection{Current pull}
With the ladder-consistent budget $\Om = 0.3096$ and the open sign, the present-epoch prediction is $\Ok^{\mathrm{obs}} =
+0.00150$ (L0) or $+0.00143$ (L1). The most recent curvature-sensitive analysis of DESI DR2 +
Planck \cite{ChenZaldarriaga2025} reports $\Ok^{\mathrm{obs}} = +0.0023 \pm 0.0011$, a mild
($2.1\sigma$) preference for an open universe. The framework pull is
$(0.0023-0.00150)/0.0011 \approx 0.72\sigma$ (L0), $0.79\sigma$ (L1) --- fully consistent, and
on the open side. A preliminary MCMC run imposing the $\Lb$-constraint over Planck plik + DESI
DR2 returns $\Ok^{\mathrm{obs}} = +0.00152$ with $\Delta\chi^2 = -10.77$ versus flat
$\Lambda$CDM and a Hubble constant $H_0 = 68.6$ shifted modestly toward the local-distance-ladder
value. We flag the chains, priors, and likelihood configuration as not yet released in
publication-grade form, and we therefore do \emph{not} treat this run as evidence for the
framework; a full, reproducible Cobaya analysis is forthcoming.

\subsection{Forecasts and discrimination}
\begin{table}[h]
\centering
\renewcommand{\arraystretch}{1.3}
\begin{tabular}{@{}l l l l@{}}
\toprule
Survey (epoch) & $\sigma_K$ & Detection of central value & L0-vs-L1 discrimination \\
\midrule
DESI DR3 ($\sim$2027) & $\sim 3\times 10^{-4}$ & $\sim 5\sigma$ & $\sim 0.24\sigma$ \\
Euclid DR1 & $\sim 1.5\times 10^{-4}$ & --- & $\sim 0.48\sigma$ \\
Euclid + CMB-S4 (2030+) & $\sim 5\times 10^{-5}$ & --- & $\sim 1.4\sigma$ \\
\bottomrule
\end{tabular}
\caption{Forecast sensitivities. DESI DR3 will test the central relation at $\sim 5\sigma$; the
L0/L1 refinement requires the combined Euclid + CMB-S4 sensitivity.}
\label{tab:forecast}
\end{table}

\section{Numerical MCMC details}
\label{sec:mcmc}
\emph{Preliminary internal results: chains, priors, and the likelihood configuration are not yet
released in publication-grade form. The numbers below are indicative only and are not used to
support the framework's claims; a complete, reproducible Cobaya chain set will be released
separately.} A Planck plik $+$ DESI DR2 BAO chain was run with Cobaya
(clik likelihoods), eight walkers, $5\times10^5$ steps. The constraint
$\Lb = 8\pi^2/\ag = 458.4$ is imposed as a derived relation $\OL = \Lb\,|\Ok|$ over a free $\Ok$,
with $\Om + \OL + \Ok = 1$. Sampled: $\Omega_b h^2$, $\Omega_c h^2$, $\theta_{\mathrm{MC}}$,
$\tau$, $n_s$, $\ln(10^{10}A_s)$, $\Ok$. Relative to flat $\Lambda$CDM the chain achieves
$\Delta\chi^2_{\mathrm{CMB}} = -6.4$ and $\Delta\chi^2_{\mathrm{BAO}} = -4.4$
($\Delta\chi^2_{\mathrm{tot}} = -10.77$), with $\Ok = +0.00152$ and $H_0 = 68.616$, modestly
shifted toward the local-distance-ladder value. Two independent likelihood implementations agree
within $\Delta\chi^2 < 1$.

\section{Consolidated status and falsifiability}
\label{sec:status}

\subsection{The epistemic ledger}
\begin{longtable}{@{}p{6.8cm} l p{4.4cm}@{}}
\caption{Consolidated status of the framework's components.}\label{tab:ledger}\\
\toprule Component & Status & Basis \\ \midrule \endfirsthead
\toprule Component & Status & Basis \\ \midrule \endhead
$Q$-charge $\int_{D^4} Q_4 = 8\pi^2$ & proven & Chang--Yang/Branson; M1-Q \\
Spectral results R1--R12 & proven & spectral computation; \cite{Vassilevich2003} \\
$\zeta(0;\mathrm{Maxwell};S^4) = -31/45$ (R1) & proven & two derivations \\
BFK identity (both bookkeepings) & proven & \cite{BFK1992}; M1-$\zeta$ \\
Hemisphere quotient $\ag/(8\pi^2)$ (L0) & proven & topological identity (Sec.~\ref{sec:main}) \\
Cap-saddle $\dtheta^*_{\mathrm{exact}} = 0.16391$ & proven & exact saddle Eq.~\eqref{eq:saddle} \\
Cap-saddle stability (all sectors) & proven (analytic) & domain monotonicity + closed-$S^4$ spectrum; M1-H \\
Boundary refinement (BFK odd-function) & proven & rigorous symmetry (Sec.~\ref{sec:boundary}) \\
Uniqueness of the invariant & proven (within criteria) & Paneitz zero-mode + (i)--(iv); full minimality open \\
KP sequestering (magnitude problem) & derived within CSG & Built-in via St\"uckelberg compensator (proof sketch); full all-orders rigor pending \\
KP self-consistency (value, structural) & numerically established & Proposition~\ref{thm:kp} (sign $G'<0$ analytic, $|G'|<2$ bound numerical); M1-$\Lambda$ \\
Three-loop bound $\approx 4.3\times 10^{-14}$ & proven & FORM; M1-3L \\
$\ag = 31/180$ & derived & Gilkey $a_4$ on $S^4$ (Sec.~\ref{sec:heat}); scalar check $1/360$ \\
$c_1 = 2\zeta(3)$ & literature; cancels from ratio & Mottola--Vaulin; convergence only, not in prediction \\
Absolute $\OL$ (KP volume) & universality structure only; scale open & $\Lambda^*$ universal in $\varphi_0$/$\Omega_m$ (proven) but amplitude-set ($\approx-0.10$ for $A_{\mathrm{CW}}$, not $-0.69$); $\varphi_0$ free; the $\langle R\rangle/4$ sharpness proxy was a floored-geometry artefact (retracted) \\
Budget ladder $\OL = 4\ag$, $\Ok = \ag^2/2\pi^2$ & conjecture (sharpened) & ladder algebra exact ($\Ok\cdot32\pi^2 = \OL^2$, $32\pi^2 = \int_{D^4}E_4$), Planck match $<1\sigma$; lift prefactor now \emph{computed} $=A/G$ (Bekenstein--Hawking horizon dof $=4S_{dS}$), residual $=$ CKN saturation P5 (named, not proven); Sec.~\ref{sec:ladder} \\
P5 localization (in-in) & localized, not proven & Schwinger--Keldysh: anomaly Adler--Bardeen exact; lift $=$ Cohen--Kaplan--Nelson bound; Schwarzschild part follows from gravity (fixes $10^{-122}$), saturation $=$ P5 \\
Scope (ratio vs absolute) & ratio holography-free & $\Ok/\OL=\ag/8\pi^2$ scale-invariant ($H,\Mpl$ cancel) $\Rightarrow$ Euclid-falsifiable core needs no P5; absolute $\OL$ alone carries the saturation \\
$\OL = 0.704$ (halo-averaged) & derived, conditional & Lombriser 2019; assumes $y(t_0)=\tfrac12$ \\
Absolute $\OL$ (causal closure) & geometric, prior-free & apparent-horizon hypersurface forced ($rH{=}1$); form not forced (KP/thermo fail, closures spread $0.49$--$0.96$) $\to$ $O(1)$ band \\
A5a (topological ratio $\ag/8\pi^2$) & theorem & Cauchy + Chern--Weil + APS; $|\epsilon_\partial|\le 1.5\%$ \\
A5b (Wick-rotation kinematics) & derived & saddle-point, conditional on B1 (established, reduced) \\
A5c (reference-state choice) & specification & the only genuinely chosen element; data-anchored \\
A2 (IR fixed point) & numerically established & contraction $L=|1+G'|<1$ requires $-2<G'<0$: sign $G'<0$ analytic, $|G'|<2$ bound numerical \\
A2 as an RG-flow attractor & retracted & not required; topology replaces it (M1-FRG) \\
Curvature sign $\Ok > 0$ (open) & forced & unique cap saddle $\dtheta^*>0$ for $\ag>0$; weighting-independent \\
Euclidean fluctuation stability & resolved (GHP) & conformal mode rotated; transverse modes positive (B1$'$) \\
Euclidean vs Lorentzian definition & open (interpretational) & FLT debate; not CSG--KP-specific \\
B1 (cap dominance, reduced action) & established (analytic) & $n_{\mathrm{complex}}=0$; no Stokes line; gap $>7\times 10^5$ \\
B1$'$ (cap $+$ scalar embedding) & complete (analytic) & degree-$8$ polynomial reduction, both $\varphi$-branches; Picard--Lefschetz $\Rightarrow n_{\mathrm{complex}}=0$ \\
Full no-boundary path integral & assembled (semiclassical) & $1$-loop det $=-4\ag$; Type-A Adler--Bardeen exact \\
All-orders stability of $R=\ag/(8\pi^2)$ & proven (within CSG) & no graviton loops; ratio of $\hbar$-indep.\ invariants; abs.\ $\varphi^4$ Borel decoupled \\
\bottomrule
\end{longtable}
\noindent
The observable prediction reduces, after this decomposition, to a single genuinely chosen
element: the reference-state specification A5c (the canonical $\Lambda$-dominated present epoch),
anchored by current cosmological data. Its Euclidean half A5a is a theorem (Cauchy + Chern--Weil
+ APS) and the Wick-rotation kinematics A5b are derived in saddle-point given B1; the
constancy of the ratio in redshift is \emph{topological} (Komargodski--Schwimmer invariance of an
integrated anomaly), so the retracted RG-attractor reading of A2 does not bear on it --- and A2
itself is \emph{numerically established} in its physically relevant form, the iterative fixed point of the
KP self-consistency map (M5 Banach contraction; sign $G'<0$ analytic, $|G'|<2$ bound numerical). The cap-saddle dominance B1 is analytically established on the reduced cap-action, and the two-dimensional cap-plus-scalar embedding $B1'$ is complete ($n_{\mathrm{complex}}=0$ by polynomial reduction, Section~\ref{sec:found}); the embedding into the full infinite-dimensional no-boundary path integral over all inhomogeneous matter and metric modes is assembled semiclassically and is one-loop exact for the Type-A prediction (Section~\ref{sec:found}); the prediction is all-orders stable as a ratio of $\hbar$-independent invariants (no graviton loops in CSG, Adler--Bardeen-exact Type-A, topological $8\pi^2$; Section~\ref{sec:found}); the one residue is the absolute Borel summability of the $\varphi^4$ matter series, which cancels from the ratio and does not bear on the prediction.

\subsection{What the framework resolves}
The three faces of the cosmological constant problem are addressed as follows. The
\emph{magnitude} is handled in two stages: the bare radiative vacuum energy is removed by
sequestering (derived within CSG, Section~\ref{sec:sequester}), while the magnitude of the anomaly-induced curvature density
rests on the holographic accumulation postulate P5 (Section~\ref{sec:psf}), which carries the
$122$-order hierarchy and is not derived from first principles. The \emph{value} is fixed
structurally by Proposition~\ref{thm:kp}: $\Lambda$ is not free, but the unique fixed point of the
averaging map; its absolute scale is a calibration at the homogeneous level. The
\emph{coincidence} is addressed by halo-weighted averaging, which returns $\OL = 0.704$
conditional on the uniform-prior assumption. The \emph{ratio}~\eqref{eq:central} --- parameter-free
given the identification A5, its topological core A5a a theorem --- links the residual
cosmological constant to the spatial curvature and supplies a single sharp test.

\paragraph{Two levels of claim.} The results separate onto two levels; conflating them is the main source of overstatement. \emph{The dimensionless prediction} --- the ratio
$\Ok/\OL = \ag/(8\pi^2)$ and the open sign $\Ok>0$ --- is theorem-level up to one data-anchored
convention. Each ingredient is proven or derived: $\ag = 31/180$ (Gilkey $a_4$), $\CQ = 8\pi^2$
(Chang--Yang), the hemisphere quotient, the cap saddle and its \emph{analytic} stability (domain
monotonicity), the contraction A2 ($L = |1+G'| < 1$; sign $G'<0$ analytic, $|G'|$ bound numerical), the dominance B1 and the complete
embedding B1$'$, the topological half A5a (theorem) and the kinematic half A5b (derived given
B1); uniqueness holds within criteria (i)--(iv). The only chosen element is the
reference epoch A5c, fixing ``today'' for a quantity scaling as $(1+z)^2$ --- a calibration
convention, not a postulate.

\emph{The absolute magnitude} --- why the curvature density has its
observed size and why $\OL \approx 0.685$ --- is a logically separate question the
dimensionless prediction does not answer. Two former ingredients have moved off the open list:
sequestering is \emph{derived within CSG} (the Built-in Sequestering
proposition, Section~\ref{sec:sequester}: the Weyl $R^2$ term supplies the
Stückelberg compensator realizing the KP constraint --- a proof sketch, full all-orders
rigor pending), and $c_1 = 2\zeta(3)$ \emph{cancels from the ratio}. What remains not derived from first principles is then narrower: the
holographic accumulation P5 (carrying the $122$-order hierarchy, its coherence
enforced by K1--K3 but not strictly derived, Section~\ref{sec:psf}); and the absolute value $\OL$
itself. The latter is subtler than a flat ``not derivable'': the self-consistency route yields a
\emph{universal} fixed point $\Lambda^*$ (independent of $\varphi_0$ and $\Omega_m$) but only
amplitude-set ($\approx-0.10$ for the CW amplitude, not the $-0.69$ once read off a
$\langle R\rangle/4$ proxy that proves a floored-geometry artefact), so $\varphi_0$ --- hence
$\OL = (V(\varphi_0)-\Lambda^*)/3$ --- stays free.
The \emph{anthropic} measure problem is categorical but \emph{avoidable}: a causal variational closure needs no measure, and on
it the boundary hypersurface is \emph{forced} (the apparent horizon, $\theta=0$, $rH=1$). What stays open is then neither measure nor boundary but the
closure \emph{functional}: KP stationarity over the
cone has no fixed point, apparent-horizon thermodynamics
give an identity or $\OL=1$, and bare dimensional closures spread across $0.49$--$0.96$ (the
$0.725$ choice not distinguished). The causal route thus
pins the hypersurface and an $O(1)$ range, not a sharp value. The framework therefore delivers a sharp
parameter-free \emph{relation}, a universal but amplitude-set
fixed point with the absolute scale open, a forced apparent-horizon
hypersurface with an $O(1)$ range for $\OL$, and --- conjecturally --- the budget ladder
($\OL = 4\ag$, matching Planck) awaiting one identification step; but \emph{not} a first-principles
value for the dark-energy density, pending either a derived closure
functional or that ladder bridge.

\subsection{Falsification criteria}
\begin{enumerate}[leftmargin=*,nosep]
\item $\Ok$ outside $\pm 3\sigma$ of $\ag\OL/(8\pi^2)$ refutes the central relation.
\item A robust detection of a \emph{closed} universe ($\Ok < 0$) refutes the sign prediction.
\item The constant object is the \emph{action-parameter} ratio $\mathcal{R}=3|K|/\Lambda$; the
density-parameter ratio redshifts as $\Ok(z)/\OL(z)=\mathcal{R}/a^2$ and is therefore \emph{not}
epoch-independent. The sharp test is the present-epoch value $\Ok^{(0)}/\OL^{(0)}=\mathcal{R}$
(extracted as a derived parameter with $a_0=1$): a measured $\Ok^{(0)}/\OL^{(0)}$ outside
$\pm3\sigma$ of $\ag/(8\pi^2)$ refutes the relation.
\end{enumerate}
DESI DR3 will test the first and third within reach; the present $\ge 2\sigma$ preference for an
open universe (DESI\,BAO${}+{}$CMB${}+{}$SNe) already constrains the second. The sign evidence is
\emph{not} yet robust across datasets, however: while the BAO-combined fits favour
$\Ok>0$, Planck temperature data alone pull toward a \emph{closed} universe
($\Ok\approx-0.04$ at ${\sim}2$--$3\sigma$), an effect entangled with the CMB lensing-amplitude
anomaly. The theoretical sign is forced ($\ag>0$), but its observational confirmation remains
data-combination dependent until DR3.

\subsection{Robustness: objections forced, cleared, or genuinely open}
A falsification attempt raised seven objections. Five are forced by the
structure or cleared by results elsewhere in the pipeline; one (DESI
$w$-dynamics) is a genuine open empirical threat we do \emph{not} claim to clear; the
seventh is the single postulate isolated above. We collect them so the one unresolved item is stated as such, not glossed over.
\begin{enumerate}[leftmargin=*,nosep]
\item \emph{The open sign is a free choice.} It is not: the reduced cap action has saddle equation
$\ag\cos^4\delta = \sin\delta$, whose only real root is $\dtheta^*=0.164>0$; a closed universe
would require $\ag<0$, so the sign $\Ok>0$ is forced by $\ag=31/180>0$.
\item \emph{The factor four in $\OL=4\ag$ is numerology.} It is the spectral ratio
$\zeta(0;\mathrm{Maxwell};S^4)/\ag = (-31/45)/(-31/180)=4$, fixed by the BFK completeness identity
$2\zeta(0;D^4)-\zeta(0;\mathrm{DtN})=-4\ag$; $\ag$ is the Gilkey coefficient and $\pi^3\approx 31$ is
flagged as coincidence, not used.
\item \emph{No dynamics produces the value, so it is unfounded.} The no-go is a feature:
$\Ok/\OL=\ag/8\pi^2$ is a topological ratio protected by conformal invariance, characterising the
geometry at all epochs, not an initial condition propagating through Friedmann evolution. Deriving
it dynamically is a category error, and the absence of any dynamical route is consistent with the
topological reading rather than evidence against it.
\item \emph{The Maxwell entanglement coefficient is $-16/45$, not $-31/45$.} The $1/3$ gap is the
gauge edge-mode contribution (Casini--Huerta--Mag\'an--Pontello), and the DtN determinant measures it
exactly: $\zeta(0;\mathrm{DtN};\mathrm{Maxwell};S^3)=3=\dim S^3$ counts the boundary polarizations,
and through the sphere factorization $Z_{S^4}=Z_{\mathrm{bulk}}/Z_{\mathrm{edge}}$ the DtN
determinant is $Z_{\mathrm{edge}}$. The discrepancy is encoded, not missed, and is itself an
independent lattice-testable prediction.
\item \emph{Why only the photon; the massless graviton should contribute.} The graviton is
non-conformal (the Lichnerowicz operator's Riemann term breaks Weyl invariance, so the conformal
result does not apply and it drops out for the same reason massive fields do), it is the geometry
being determined (its inclusion is circular), and on $D^4$ only $\chi$ survives ($\sigma=0$). The
photon is the unique massless conformal non-gravitational field.
\item \emph{DESI prefers dynamical dark energy at $>3\sigma$, refuting $w=-1$.} This is the one
objection that is \emph{not} cleared. The hope that DESI's $w_0w_a$ preference (obtained under
forced $\Ok=0$) is an artefact of the $\Ok$--$w_0$ degeneracy fails quantitatively: opening to
$\Ok=+1.5\times10^{-3}$ shifts the effective $w_0$ by only ${\sim}2\times10^{-3}$, whereas DESI
requires $0.25$--$0.6$---a factor ${\sim}100$--$300$ too small, and the running-vacuum strength
$\nu=\ag/8\pi^2\approx 2.2\times10^{-3}$ is equally small. The framework therefore carries
\emph{two} independent tensions with DESI: curvature (${\sim}1.8\sigma$, tolerable) and $w$
(${\sim}4\sigma$ if real), and the second cannot be absorbed by the first. Honest status: either
DESI's $w$-signal weakens with data, or $w\neq-1$ is real and the constant $\OL=4\ag$ requires an
extension. This is a genuine open empirical threat, decided by DESI DR3 and Euclid, not by the
present analysis.
\end{enumerate}
The honest boundary is therefore exactly \emph{one conceptual postulate} --- the anomaly--curvature identification $\mathrm{A5}=\text{(iii)}=\mathrm{P5}$ --- \emph{together with one sharp empirical threat} ($w\neq-1$), both decided by observation.

\section*{Acknowledgement of conditional status}
\addcontentsline{toc}{section}{Acknowledgement of conditional status}
This paper separates what is proven from what is assumed. The Euclidean conformal and
topological core --- the $Q$-charge, the twelve spectral results, the BFK identity, the
hemisphere quotient, the cap-saddle and its stability, the boundary refinement, the uniqueness
statement, and the structural content of the self-consistency theorem --- is established at
theorem level and independently recomputed (M1). The photon anomaly coefficient $\ag = 31/180$ is
derived from the Gilkey heat-kernel coefficient (Section~\ref{sec:heat}); the two-loop
normalization $c_1 = 2\zeta(3)$ enters only the convergence bound and cancels from the predicted
ratio. Sequestering is derived within CSG (the Built-in Sequestering proposition: the Weyl
$R^2$ term supplies the Stückelberg compensator realizing the KP constraint --- a proof sketch,
full all-orders rigor pending), not an external add-on. The absolute dark-energy
scale is a calibration, becoming a conditional derivation only with the structure-formation
assumption $y(t_0)=\tfrac12$. The KP self-consistency fixed point A2 is numerically established (M5 Banach
contraction; sign analytic, $|G'|$ bound numerical); cap-saddle dominance B1 is analytically established on the reduced cap-action; and
the Lorentzian identification A5 decomposes into a theorem (A5a), a B1-conditional kinematic
derivation (A5b), and a single data-anchored choice (A5c). The open curvature sign is forced by
$\ag>0$ via the unique cap saddle. What remains genuinely open is: the holographic accumulation
postulate P5 (which carries the $122$-order magnitude hierarchy and lacks a Schwinger--Keldysh
derivation of $N_{\mathrm{eff}} = \Mpl^2/H^2$); the choice A5c; the absolute Borel summability of
the $\varphi^4$ matter series (the prediction itself being all-orders stable as a ratio of
$\hbar$-independent invariants, Section~\ref{sec:found}); and the Lorentzian
fluctuation stability (the Feldbrugge--Lehners--Turok question) --- the last of these is where
observation, not further algebra, will decide.

\section{Code and data availability}
\label{sec:code}
A self-contained Python package reproduces every numerical and symbolic statement of this
paper. The headline constants and the central prediction are defined once in
\texttt{csg\_kp\_core.py} (the single source for $\ag = 31/180$, $\CQ = 8\pi^2$, the matter
fraction $\Om = 0.3096$, and the derived $\OL$, $\Ok$); all other modules import from it. The
full verification harness is \texttt{run\_all.py} ($37/37$ modules), with the dedicated
falsification and audit harnesses \texttt{foundations/falsification\_suite.py} ($16/16$ checks)
and \texttt{foundations/axiom\_status\_audit.py}. Individual claims map to modules as follows.
\begin{itemize}[leftmargin=*,nosep]
\item $\ag = 31/180$ from the Gilkey heat-kernel coefficient (independent of the core, to verify
it): \texttt{foundations/a\_gamma\_derivation.py}.
\item Universality of the Banach fixed point $\Lambda^*$ and the Coleman--Weinberg iteration:
\texttt{foundations/cw\_banach\_iteration.py}.
\item Coherent spectral accumulation lifting the local anomaly to the curvature scale:
\texttt{foundations/spectral\_coherence.py}.
\item The geometric budget ladder $\OL = 4\ag$ from $\ag$ and $\chi = 2$:
\texttt{foundations/budget\_ladder.py}.
\item The $A/G$ lift prefactor (Bekenstein--Hawking horizon degrees of freedom):
\texttt{foundations/action\_density\_bridge.py}.
\item The Schwinger--Keldysh localization of the residual assumption P5 (CKN saturation):
\texttt{foundations/schwinger\_keldysh\_p5.py}.
\item The scope confinement (conjecture constrains the absolute value, not the ratio):
\texttt{foundations/ratio\_vs\_absolute\_scope.py}.
\item The no-go audit for the absolute $\OL = 4\ag$ value:
\texttt{foundations/absolute\_value\_audit.py}.
\item The three residual foundational points (sign, $B1'$, reference epoch):
\texttt{foundations/close\_open\_points.py}.
\end{itemize}

\begin{thebibliography}{99}
\bibitem{Bekenstein1972} J.~D.~Bekenstein, \textit{Lett.\ Nuovo Cimento} \textbf{4}, 737 (1972).
\bibitem{Strominger1996} A.~Strominger and C.~Vafa, \textit{Phys.\ Lett.\ B} \textbf{379}, 99 (1996); arXiv:hep-th/9601029.
\bibitem{DS1993} S.~Deser and A.~Schwimmer, \textit{Phys.\ Lett.\ B} \textbf{309}, 279 (1993); arXiv:hep-th/9302047.
\bibitem{Duff1994} M.~J.~Duff, \textit{Class.\ Quantum Grav.} \textbf{11}, 1387 (1994).
\bibitem{KomargodskiSchwimmer2011} Z.~Komargodski and A.~Schwimmer, \textit{JHEP} \textbf{12}, 099 (2011); arXiv:1107.3987.
\bibitem{Birrell1982} N.~D.~Birrell and P.~C.~W.~Davies, \textit{Quantum Fields in Curved Space} (Cambridge, 1982).
\bibitem{ChristensenDuff1980} S.~M.~Christensen and M.~J.~Duff, \textit{Nucl.\ Phys.\ B} \textbf{170}, 480 (1980).
\bibitem{ChangYang1995} S.-Y.~A.~Chang and P.~C.~Yang, \textit{Annals of Math.} \textbf{142}, 171 (1995).
\bibitem{Branson1995} T.~P.~Branson, \textit{Trans.\ Amer.\ Math.\ Soc.} \textbf{347}, 3671 (1995).
\bibitem{BFK1992} D.~Burghelea, L.~Friedlander, T.~Kappeler, \textit{J.\ Funct.\ Anal.} \textbf{107}, 34 (1992).
\bibitem{Vassilevich2003} D.~V.~Vassilevich, \textit{Phys.\ Rept.} \textbf{388}, 279 (2003).
\bibitem{KP2014} N.~Kaloper and A.~Padilla, \textit{Phys.\ Rev.\ Lett.} \textbf{112}, 091304 (2014); arXiv:1309.6562.
\bibitem{ShawBarrow2011} D.~J.~Shaw and J.~D.~Barrow, \textit{Phys.\ Rev.\ D} \textbf{83}, 043518 (2011); arXiv:1010.4262.
\bibitem{BD2015} I.~Ben-Dayan et al., \textit{JCAP} \textbf{05}, 002 (2016); arXiv:1507.04158.
\bibitem{Riegert1984} R.~J.~Riegert, \textit{Phys.\ Lett.\ B} \textbf{134}, 56 (1984).
\bibitem{BGT1995} M.~Bucher, A.~S.~Goldhaber, N.~Turok, \textit{Phys.\ Rev.\ D} \textbf{52}, 3314 (1995); arXiv:astro-ph/9411046.
\bibitem{HawkingTurok1998} S.~W.~Hawking and N.~Turok, \textit{Phys.\ Lett.\ B} \textbf{425}, 25 (1998).
\bibitem{Feldbrugge2017} J.~Feldbrugge, J.-L.~Lehners, N.~Turok, \textit{Phys.\ Rev.\ Lett.} \textbf{119}, 171301 (2017).
\bibitem{Barbour2014} J.~Barbour, T.~Koslowski, F.~Mercati, \textit{Phys.\ Rev.\ Lett.} \textbf{113}, 181101 (2014).
\bibitem{Lombriser2019} L.~Lombriser, \textit{Phys.\ Lett.\ B} \textbf{797}, 134804 (2019); arXiv:1901.08588.
\bibitem{MottolaVaulin2006} E.~Mottola and R.~Vaulin, \textit{Phys.\ Rev.\ D} \textbf{74}, 064004 (2006); arXiv:gr-qc/0604051.
\bibitem{Planck2018} Planck Collaboration, \textit{A\&A} \textbf{641}, A6 (2020); arXiv:1807.06209.
\bibitem{ChenZaldarriaga2025} J.~Chen and M.~Zaldarriaga, \textit{JCAP} \textbf{08}, 027 (2025); arXiv:2505.00659.
\bibitem{Bohr1928} N.~Bohr, \textit{Nature} \textbf{121}, 580 (1928).
\bibitem{ReehSchlieder1961} H.~Reeh and S.~Schlieder, \textit{Nuovo Cim.} \textbf{22}, 1051 (1961).
\bibitem{Hensen2015} B.~Hensen et al., \textit{Nature} \textbf{526}, 682 (2015); arXiv:1508.05949.
\bibitem{Giustina2015} M.~Giustina et al., \textit{Phys.\ Rev.\ Lett.} \textbf{115}, 250401 (2015); arXiv:1511.03190.
\bibitem{Shalm2015} L.~K.~Shalm et al., \textit{Phys.\ Rev.\ Lett.} \textbf{115}, 250402 (2015); arXiv:1511.03189.
\bibitem{CohenKaplanNelson1999} A.~G.~Cohen, D.~B.~Kaplan and A.~E.~Nelson, \textit{Phys.\ Rev.\ Lett.} \textbf{82}, 4971 (1999); arXiv:hep-th/9803132.
\bibitem{Li2004} M.~Li, \textit{Phys.\ Lett.\ B} \textbf{603}, 1 (2004); arXiv:hep-th/0403127.
\bibitem{Hsu2004} S.~D.~H.~Hsu, \textit{Phys.\ Lett.\ B} \textbf{594}, 13 (2004); arXiv:hep-th/0403052.
\bibitem{PavonZimdahl2005} D.~Pav\'on and W.~Zimdahl, \textit{Phys.\ Lett.\ B} \textbf{628}, 206 (2005); arXiv:gr-qc/0505020.
\bibitem{Mottola2010} E.~Mottola, arXiv:1006.3567 (2010).
\bibitem{Mottola2022} E.~Mottola, \textit{JHEP} \textbf{11}, 037 (2022); arXiv:2205.04703.
\bibitem{Jacobson2016} T.~Jacobson, \textit{Phys.\ Rev.\ Lett.} \textbf{116}, 201101 (2016); arXiv:1505.04753.
\bibitem{RyuTakayanagi2006} S.~Ryu and T.~Takayanagi, \textit{Phys.\ Rev.\ Lett.} \textbf{96}, 181602 (2006); arXiv:hep-th/0603001.
\bibitem{JacobsonVisser2019} T.~Jacobson and M.~R.~Visser, \textit{SciPost Phys.} \textbf{7}, 079 (2019); arXiv:1812.01596.
\bibitem{Verlinde2017} E.~P.~Verlinde, \textit{SciPost Phys.} \textbf{2}, 016 (2017); arXiv:1611.02269.
\bibitem{PadmanabhanCosMIn2013} H.~Padmanabhan and T.~Padmanabhan, \textit{Int.\ J.\ Mod.\ Phys.\ D} \textbf{22}, 1342001 (2013); arXiv:1302.3226.
\bibitem{BilicEtAl2007} N.~Bili\'c, B.~Guberina, R.~Horvat, H.~Nikoli\'c and H.~\v{S}tefan\v{c}i\'c, \textit{Phys.\ Lett.\ B} \textbf{657}, 232 (2007); arXiv:0707.3830.
\bibitem{SolaPeracaula2021} J.~Sol\`a Peracaula, \textit{Int.\ J.\ Mod.\ Phys.\ A} \textbf{36}, 2150239 (2021); arXiv:2109.12086.
\end{thebibliography}

\end{document}
